\documentclass[11pt]{article}
\usepackage[T1]{fontenc}
\usepackage[utf8]{inputenc}
\usepackage{lmodern}
\usepackage[a4paper,margin=0.88in]{geometry}
\usepackage{amsmath,amssymb,amsthm,mathtools}
\usepackage{microtype}
\usepackage{booktabs,tabularx,longtable,array}
\usepackage{enumitem}
\usepackage{xcolor}
\usepackage[hidelinks]{hyperref}
\makeatletter
\g@addto@macro\UrlBreaks{\do\_\do\-}
\makeatother
\usepackage[nameinlink,capitalize]{cleveref}
\usepackage{fancyhdr}
\usepackage{lastpage}

\IfFileExists{buildmeta.tex}{\input{buildmeta.tex}}{%
  \newcommand{\BuildCommit}{local-audit}%
  \newcommand{\BuildDate}{28 August 2026}%
}

\definecolor{deepblue}{RGB}{25,55,92}
\definecolor{softblue}{RGB}{235,242,249}
\definecolor{softgray}{RGB}{245,245,245}
\definecolor{softgreen}{RGB}{237,247,239}
\definecolor{softred}{RGB}{252,239,239}
\definecolor{softamber}{RGB}{255,247,230}

\newtheorem{theorem}{Theorem}[section]
\newtheorem{proposition}[theorem]{Proposition}
\newtheorem{lemma}[theorem]{Lemma}
\newtheorem{corollary}[theorem]{Corollary}
\theoremstyle{definition}
\newtheorem{definition}[theorem]{Definition}
\newtheorem{remark}[theorem]{Remark}

\newcommand{\N}{\mathbb N}
\newcommand{\Q}{\mathbb Q}
\newcommand{\R}{\mathbb R}
\newcommand{\C}{\mathbb C}
\newcommand{\1}{\mathbf 1}
\newcommand{\E}{\mathcal E}
\newcommand{\B}{\mathcal B}
\newcommand{\WPS}{\operatorname{WindowPairSupply}}
\newcommand{\status}[1]{\textsf{#1}}
\newcommand{\lean}{\textsc{Lean}}

\pagestyle{fancy}
\fancyhf{}
\lhead{Erd\H{o}s \#287 -- V16.3 Theory Manual R4}
\rhead{28 August 2026}
\cfoot{\thepage\ of \pageref{LastPage}}
\setlength{\headheight}{14pt}
\setlength{\emergencystretch}{2em}

\title{\vspace{-1.2cm}\textbf{A Machine-Checked Finite Reduction for Erd\H{o}s Problem \#287\\
and a Source-Explicit Large-$M$ Criterion}\\[1ex]
\large Clean Public Review + Self-Contained Gate Theory Manual\\[0.3ex]\normalsize Version 16.3 -- Prepublication Candidate R4 (not yet published)}
\author{Research draft with Lean/Aristotle-assisted formal verification}
\date{28 August 2026}

\begin{document}
\maketitle

\begin{center}
\fcolorbox{deepblue}{softblue}{\begin{minipage}{0.92\textwidth}
\textbf{Status.} Erd\H{o}s Problem \#287 remains open.  The unconditional public
Lean result excludes every exact counterexample with maximum denominator at most
$4\times10^9$, and the large-$M$ conclusion is compiled from the explicit
\texttt{WindowPairSupply} interface.  A later V15 Aristotle replay, supplied with this audit,
reports a successful $8118$-job kernel build of the source-minimal identity
$\Lambda=\mu*\log$, its affine source, the determinant-one line, the balanced-seven
coefficient $-20$, and the squarefree labelled polarization.  Public repository
\texttt{main} already contains individual V15 source modules, but its aggregate \texttt{Main.lean}/status replay is not yet synchronized; this note therefore distinguishes
\emph{Lean-proved in the latest supplied replay} from \emph{public-main synchronized}.
The latest supplied V16 factorial Aristotle replay reports the factorial Euler seven-box
polarization, its repeated-prime fibre identity, the local Euler/generalized-von-Mangoldt
algebra, expected-term linearity and the Pascadi parameter ledger as kernel-checked in an
$8123$-job build.  Public-main synchronization remains a separate provenance question.  The first exact balanced-seven analytic open is
\path{AFFINE287-FACTORIAL-OMEGA7-SIGNED-ENDPOINT45}; the physical comparison match,
remaining packets, large-$M$ supply, and Problem \#287 remain open.  This R4 candidate
has not been pushed to the stable public PDF.
\end{minipage}}
\end{center}

\begin{abstract}
Erd\H{o}s Problem \#287 asks whether every representation
\[
1=\frac1{n_1}+\cdots+\frac1{n_k},\qquad 1<n_1<\cdots<n_k,
\]
contains a consecutive gap of size at least $3$.  The present project has one unconditional
machine-checked result and one explicit conditional large-$M$ interface.  Lean excludes
all exact counterexamples with maximum denominator at most $4\times10^9$, while an
effective eventual \texttt{WindowPairSupply} theorem would close the remaining range.

For the analytic programme we use the source-minimal identity
\[
\Lambda(2mn+s)=\sum_{qr=2mn+s}\mu(q)\log r,\qquad s=\pm1,
\]
together with its exact small-$q$/small-$r$/hard partition and determinant-one affine line.
A supplied V15 Aristotle replay reports these identities, the balanced-seven coefficient
$-20$, and the earlier squarefree polarization as kernel-checked.  The current balanced-seven
representation is sharper: for
\[
a_{\mathbf z}(p)=\frac17\sum_{i=1}^7z_i\omega_i(p)
\]
we introduce the factorial Euler multiplicative function
\[
F_{\mathbf z}(p^e)=\frac{a_{\mathbf z}(p)^e}{e!}.
\]
Its degree-$(1,\ldots,1)$ torus coefficient recovers the ordered seven labelled prime-box
source \emph{including repeated primes exactly}; hence the previous
$X^{6/7+o(1)}$ collision remainder is no longer part of the controlling formula.
Polarization also commutes with every linear expected-term operator.  The remaining
balanced-seven inputs are therefore (i) one signed factorial endpoint distribution estimate
at $(Q,y)=(X^{3/5},X^{1/7})$, and (ii) identification of the factorial expected term with
the physical comparison/main term.

The endpoint estimate required here is strictly narrower than a Bombieri--Vinogradov theorem
for all class-$\mathcal C$ functions, an absolute modulus average, a theorem uniform in
$\mathbf z$, or a general smooth-number theorem.  A direct audit of Pascadi's current
arXiv v2 proof parameters is nonclosing for this particular route: the proof architecture
uses $Q\le x^{5/8-100\eta}$ together with $y\le x^\eta$; setting
$Q=x^{3/5}$ forces $\eta\le1/4000$, whereas $y=x^{1/7}$ requires
$\eta\ge1/7$.  This is a limitation of that written proof route, not a theorem excluding
future extensions.  The paper also states the Gate 0--2 theory itself in theorem form, so that each source, target, pin and non-implication can be audited without the historical dossier.  No analytic endpoint estimate, comparison identification, eventual
large-$M$ supply, or proof of Erd\H{o}s \#287 is claimed.
\end{abstract}

\section*{Status at a glance}
\addcontentsline{toc}{section}{Status at a glance}
\begin{center}
\footnotesize
\renewcommand{\arraystretch}{1.18}
\begin{tabularx}{\textwidth}{>{\raggedright\arraybackslash}p{0.29\textwidth}
>{\raggedright\arraybackslash}p{0.22\textwidth}X}
\toprule
Item & Status & Meaning \\
\midrule
Exact public counterexample predicate & \status{PUBLIC LEAN-CHECKED} & Exact set formulation and ordered public statement are connected. \\
Finite range $M\le4\times10^9$ & \status{PUBLIC LEAN-CHECKED} & A certificate chain excludes every exact counterexample in this range. \\
Window-pair closure compiler & \status{PUBLIC LEAN-CHECKED / CONDITIONAL} & Eventual effective \texttt{WindowPairSupply} plus the finite theorem implies \#287. \\
$\Lambda=\mu*\log$, affine source, determinant-one line & \status{LEAN-PROVED IN LATEST V15 REPLAY} & User-supplied V15 report records an $8118$-job kernel build; public-main synchronization is still pending. \\
Balanced-seven $-20$ and squarefree polarization & \status{LEAN-PROVED IN LATEST V15 REPLAY} & Finite/algebraic only; no analytic cancellation follows. \\
Factorial Euler seven-box polarization & \status{LEAN-PROVED IN LATEST V16 FACTORIAL REPLAY} & Exact ordered-source encoding including repeated primes; analytic endpoint remains open. \\
Expected-term/polarization linearity & \status{RESEARCH/ALGEBRAIC PASS} & Coefficient extraction commutes with every linear expected-term operator. \\
Gate 0 & \status{MIXED / NOT CLOSED} & $b.1$ and growth are controlled; uniform $b.2$, the literal Type-I source theorem and comparison/source pins remain open or external. \\
Gate 1A & \status{CONDITIONAL POSSIBLE PROVIDER} & Canonical common-weight mean-square theorem under explicit source pins; current source/fibre instantiation remains open and its necessity for \#287 is census-dependent. \\
Gate 1B & \status{OPEN COMPLEMENTARY PROVIDER} & Centered/M\"obius source algebra is banked; a source-to-convolution-BV/provider theorem remains open. \\
Ford generated-packet census & \status{OPEN SOURCE/REASSEMBLY} & Universal Type II is stronger than the proof-specific family actually generated at Ford--Maynard (7.23). \\
Gate 2 & \status{CONDITIONAL ENDGAME / NOT ACTIVATED} & Consumes Gate-0 data, a complete Type-II package, the $N_2$ replacement and compatible $\varepsilon$ to yield positive prime mass. \\
Direct Pascadi v2 proof route at $(3/5,1/7)$ & \status{AUDITED NONCLOSING} & Written parameters force $\eta\le1/4000$ but the source requires $\eta\ge1/7$. \\
Factorial signed endpoint estimate & \status{OPEN ANALYTIC} & \texttt{FACTORIAL-OMEGA7-SIGNED-ENDPOINT45} (full programme label given in Section~\ref{sec:factorial-endpoint}). \\
Physical comparison identification & \status{OPEN SOURCE} & Principal local density plus the literal exceptional-character convention must match the physical comparison term. \\
Remaining packets / FCL / eventual supply & \status{OPEN} & Balanced-seven closure alone does not finish the large-$M$ programme. \\
Erd\H{o}s Problem \#287 & \status{OPEN} & No unconditional solution claim. \\
\bottomrule
\end{tabularx}
\end{center}

\paragraph{Status vocabulary.}
\status{PUBLIC LEAN-CHECKED} means present on public repository \texttt{main}.
\status{LEAN-PROVED IN LATEST V15 REPLAY} means the supplied Aristotle V15 build report
records a successful kernel replay, while public-main synchronization remains a separate
provenance step.  \status{LEAN-PROVED IN LATEST V16 FACTORIAL REPLAY} means the exact mathematical
derivation has passed the current hostile audit but no completed Lean replay is claimed.
\status{OPEN ANALYTIC} and \status{OPEN SOURCE} denote deliberately unproved inputs.
\status{AUDITED NONCLOSING} refers only to the literal cited theorem/proof dictionary, not
to a statement that no future extension can work.

\tableofcontents

\section{The problem and the scope of this note}

The public problem asks whether, for distinct integers
\[
1<n_1<\cdots<n_k,\qquad \sum_{i=1}^k\frac1{n_i}=1,
\]
one necessarily has
\[
\max_{1\le i<k}(n_{i+1}-n_i)\ge3.
\]
It is listed as open in the Erd\H{o}s Problems archive and originates in the sources
cited there \cite{ErdosGraham1980,ErdosProblems287}.

For a finite set $A\subset\N$, write
\[
N=\min A,\qquad M=\max A.
\]
A counterexample is a finite set satisfying
\[
A\subset\{2,3,\dots\},\qquad \sum_{a\in A}\frac1a=1,
\]
and the local gap condition that every nonmaximal $a\in A$ is followed by $a+1$ or $a+2$
in $A$.  The Lean development formalizes exactly this predicate and proves the ordered
form from it.

This note is deliberately not a chronological research diary.  It records only the
current controlling finite theorem, the exact conditional closure interface, the
source-minimal algebra, the factorial Euler balanced-seven reduction, and the current open
analytic obligations.  Earlier squarefree/collision routers, Vaughan decompositions, failed
provider dictionaries, hostile audits, and superseded residuals belong to the separate technical
verification ledger and immutable V15 archive.

\section{The machine-checked finite theorem}

\begin{theorem}[Finite exclusion]
There is no exact counterexample to Erd\H{o}s Problem \#287 with
\[
3\le M\le4{,}000{,}000{,}000.
\]
Equivalently, every unit-fraction representation in the problem with maximum denominator
at most $4\times10^9$ has a consecutive gap at least $3$.
\end{theorem}

The theorem is represented in the repository by
\texttt{no\_Erdos287Counterexample\_}\\\texttt{of\_max\_le\_4e9}.  Its proof is not an exhaustive
enumeration of finite subsets.  It is a certificate compiler built from the three ideas
below.

\subsection{Holes and the adjacent-hole contradiction}

For a hypothetical counterexample, call an integer $x\in[N,M]$ a \emph{hole} when
$x\notin A$.  The gap-$\le2$ condition implies that two consecutive integers in
$[N,M]$ cannot both be holes.  This is the formal theorem
\texttt{Gap2CE.holes\_isolated}.

A second elementary fact places the top half of the interval above the minimum:
\[
N\le\left\lfloor\frac M2\right\rfloor.
\]
Thus, whenever $M\le2x$ and $x+1\le M$, the consecutive pair $x,x+1$ lies inside the
closed window where the adjacent-hole contradiction applies.

\subsection{Prime-power windows}

For $j\ge0$, let $C(j)$ denote the largest reduced numerator arising from a nonempty
partial reciprocal sum
\[
\sum_{s\in S}\frac1s,\qquad \varnothing\ne S\subseteq\{1,\dots,j\}.
\]
The formal development uses a certified upper table
\[
C_{\mathrm{val}}(0),\dots,C_{\mathrm{val}}(9)
=1,1,3,11,25,137,137,1019,2143,7129.
\]

The prime-power window theorem says, schematically, that if $p$ is prime, $e\ge1$,
$p^e\mid x$, and
\[
C\!\left(\left\lfloor\frac M{p^e}\right\rfloor\right)<p,
\]
then $x$ cannot belong to $A$.  The reason is that the top $p$-adic layer of the reciprocal
sum would force a nonzero numerator to vanish modulo $p$.

Consequently, if both $x$ and $x+1$ possess such prime-power divisors, then they are two
adjacent holes.  This gives the central finite compiler.

\begin{proposition}[Window-pair blocker]
Suppose there are primes $p_u,p_v$, positive exponents $a_u,a_v$, and an integer $x$ such
that
\[
p_u^{a_u}\mid x,\qquad p_v^{a_v}\mid x+1,
\]
\[
\left\lfloor\frac M{p_u^{a_u}}\right\rfloor,
\left\lfloor\frac M{p_v^{a_v}}\right\rfloor\le9,
\]
\[
C_{\mathrm{val}}\!\left(\left\lfloor\frac M{p_u^{a_u}}\right\rfloor\right)<p_u,
\qquad
C_{\mathrm{val}}\!\left(\left\lfloor\frac M{p_v^{a_v}}\right\rfloor\right)<p_v,
\]
and
\[
M\le2x,
\qquad x+1\le M.
\]
Then a gap-$\le2$ counterexample with maximum $M$ cannot exist.
\end{proposition}

\subsection{The interval-certificate chain}

One certificate can exclude an entire interval $L\le M\le U$.  The divisibility data are
fixed at a chosen $x$, monotonicity controls the windows as $M$ varies, and $U\le2x$
keeps the forced holes above the minimum.  The repository contains a chain of $34$
certificates whose intervals meet end-to-end and cover
\[
[3,4\times10^9].
\]
All primality side conditions are discharged in the kernel; the construction does not use
\texttt{native\_decide}.


\subsection{Sign-sensitive Sophie witnesses and the free maximum blocker}

The latest supplied finite bank sharpens the final analytic-to-finite landing and is worth stating because it
prevents small endpoint mistakes later.  A \emph{plus} Sophie witness at $M$ is a prime $q$ with
\[
 q\ge5,\qquad M<3q,\qquad 2q+1\le M,\qquad 2q+1\ \hbox{prime},
\]
equivalently $M/3<q\le (M-1)/2$ with the integer endpoints understood literally.  A \emph{minus} witness is
\[
 q\ge5,\qquad M<3q,\qquad 2q\le M,\qquad 2q-1\ \hbox{prime},
\]
so the correct minus endpoint is
\[
 \boxed{q\le\left\lfloor\frac M2\right\rfloor.}
\]
Either witness contradicts an exact counterexample.  The hypothesis $q\ge5$ is not cosmetic: for a prime it is
the same as $q>3$, which is the finite window threshold needed by the cofactor-two obstruction.

There is also a free maximum-divisor blocker.  If $q>3$ is prime, $q\mid M$ and $M<3q$, then necessarily
$M\in\{q,2q\}$; in either case the top-layer/cofactor-two exclusion contradicts $M\in A$.  Thus a counterexample
cannot have prime maximum and cannot have $M=2q$ with $q>3$ prime.  These are finite repairs, not analytic
inputs; they should be used when translating positive affine mass into an actual blocker.

\section{The exact large-$M$ closure criterion}

\begin{definition}[Window-pair supply]
For $M\in\N$, $\WPS(M)$ is the assertion that there exist
$x,p_u,a_u,p_v,a_v\in\N$ such that
\[
\begin{gathered}
p_u,p_v\text{ are prime},\qquad a_u,a_v\ge1,\\
p_u^{a_u}\mid x,\qquad p_v^{a_v}\mid x+1,\\
\left\lfloor\frac M{p_u^{a_u}}\right\rfloor\le9,
\quad C_{\mathrm{val}}\!\left(\left\lfloor\frac M{p_u^{a_u}}\right\rfloor\right)<p_u,\\
\left\lfloor\frac M{p_v^{a_v}}\right\rfloor\le9,
\quad C_{\mathrm{val}}\!\left(\left\lfloor\frac M{p_v^{a_v}}\right\rfloor\right)<p_v,\\
M\le2x,
\qquad x+1\le M.
\end{gathered}
\]
\end{definition}

In plain terms, two consecutive integers in the top half of $[1,M]$ each have a
prime-power divisor larger than roughly $M/10$, with a small local numerator condition.

\begin{theorem}[End-to-end conditional compiler]
Assume there is an explicit $M_0\le4\times10^9$ such that
\[
\WPS(M)\qquad\text{for every }M\ge M_0.
\]
Then Erd\H{o}s Problem \#287 is true.
\end{theorem}

This implication is machine-checked as
\texttt{no\_Erdos287Counterexample\_of\_closure}.  The proof splits at
$4\times10^9$: the finite theorem handles the lower range, while the window-pair blocker
handles every larger maximum.

The interface is weaker than a Sophie-Germain prime condition.  For example, if a prime
$q$ in the interval $M/3<q\le(M-1)/2$ has $2q+1$ prime, then one may take
$x=2q$.  The analogous minus case uses $x=2q-1$.  The Lean theorem
\texttt{windowPairSupply\_of\_sophieWitness} formalizes this implication.  The present
programme does not claim such Sophie witnesses for every large $M$.



\clearpage
\part{Theory and Verification Manual: Gate 0--2}
\label{part:theory-manual}

\section{What the Gate theory is trying to do}
\label{sec:theory-purpose}

This Part is written for a reader who has never seen the earlier chats, the V15 patch stack,
or the repository.  The Gate labels are not theorem substitutes: each Gate is a literal
mathematical package with an input source, a target theorem, a verification boundary, and
specified downstream consequences.

Conceptually, the programme separates four jobs.

\begin{description}[leftmargin=2.3cm,style=nextline]
\item[Gate 0.] Construct and validate the physical sequence/comparison class: prime mass,
comparison regularity, growth, Type I, and the exact expected-term convention.
\item[Gate 1.] Supply parity-breaking Type-II estimates for the actual generated packets.
\item[Gate 1A.] One canonical common-weight moving-source provider.
\item[Gate 1B.] A complementary signed/M\"obius/high-order provider that preserves the
arithmetic sign resource through completion/dispersion.
\item[Gate 2.] Consume a source-exhaustive Type-II package, the Ford arithmetic weight,
the exceptional $N_2$ sector, and one compatible $\varepsilon_\ast$ to run the
Ford--Maynard lower-bound endgame.
\end{description}

The intended dependency graph is
\[
\boxed{
\begin{array}{c}
\text{PHYSICAL \#287 SOURCE}\\
\downarrow\\
\text{GATE 0: source/comparison/Type-I validation}\\
\downarrow\\
\text{FORD GENERATED-PACKET CENSUS}\\
\downarrow\\
\begin{array}{ccc}
\text{GATE 1A}&&\text{GATE 1B}\\[-1mm]
\text{common-weight provider}&&\text{signed/M\"obius provider}
\end{array}\\
\downarrow\\
\text{SOURCE-EXHAUSTIVE TYPE-II REASSEMBLY}\\
\downarrow\\
\text{GATE 2: source-specific Ford lower-bound compiler}\\
\downarrow\\
\text{POSITIVE AFFINE MASS + EFFECTIVITY}\\
\downarrow\\
\text{EFFECTIVE LARGE-}M\text{ WITNESS}\\
\downarrow\\
\WPS(M)\\
\downarrow\\
\text{finite Lean compiler}\\
\downarrow\\
\text{ERD\H{o}S \#287.}
\end{array}}
\tag{Theory-DAG}
\]
This is a dependency graph, not a proof.

The following non-implications are load-bearing:
\[
\boxed{
\begin{gathered}
\text{Gate 1A closed}\not\Rightarrow\text{Gate 1B closed},\qquad
\text{Gate 1B closed}\not\Rightarrow\text{Gate 1A closed},\\
\text{Gate 1A + Gate 1B}\not\Rightarrow\text{Full Type II without source exhaustiveness},\\
\text{generated Type II}\not\Rightarrow\text{Gate 2 without Gate-0/$N_2$/source inputs},\\
\text{Balanced7 closed}\not\Rightarrow\text{FCL closed},\\
\text{FCL closed}\not\Rightarrow\WPS(M),\\
\text{none of the preceding statements alone proves Erd\H{o}s \#287.}
\end{gathered}}
\tag{Theory-firewall}
\]

\subsection{Verification vocabulary}

Every major object below is assigned one of the following meanings.
\begin{center}
\renewcommand{\arraystretch}{1.16}
\begin{tabularx}{\textwidth}{>{\ttfamily\raggedright\arraybackslash}p{0.26\textwidth}X}
\toprule
Label & Meaning\\
\midrule
LEAN-PROVED & Proved by a cited/current Lean kernel replay.\\
PROVED-FINITE / PROVED-ALGEBRAIC & Exact finite/combinatorial/algebraic mathematics;
may or may not also be Lean formalized.\\
EXTERNALLY-AUDITED & A sourced mathematical theorem or source audit, not formalized here.\\
CONDITIONAL-COMPILER & The implication is proved, but one or more explicit antecedents remain open.\\
SOURCE-OPEN & The physical/source identification, normalization or dictionary is missing.\\
ANALYTIC-OPEN & The source object is sufficiently fixed; a deep estimate remains unproved.\\
SOURCE-BLOCKED & The literal source object/prefactor needed even to state the final theorem is not yet printed.\\
CAPACITY-ONLY & Exponent arithmetic passes, but no analytic theorem follows.\\
FALSE / RETIRED & The literal proposed route is disproved or withdrawn.\\
SUPERSEDED & Historically valid but no longer the controlling representation.\\
\bottomrule
\end{tabularx}
\end{center}

\section{Gate 0: the physical sequence/comparison theorem}
\label{sec:gate0-theory}

\subsection{The physical \#287 sequence}

Fix a nonnegative smooth weight $W$ supported inside the analytic blocker band; the
finite blocker itself only needs the corresponding interval to lie in the top half.
Define
\[
\mathfrak S_2
=
\prod_{\ell>2}\frac{\ell(\ell-2)}{(\ell-1)^2},
\qquad
B(n)
=
\mathfrak S_2
\prod_{\substack{\ell\mid n\\\ell>2}}
\frac{\ell-1}{\ell-2},
\]
and
\[
\boxed{
a_X(n)
=
W(n/X)\bigl(\Lambda(2n-1)+\Lambda(2n+1)\bigr),
}
\]
\[
\boxed{
b_X(n)=4W(n/X)B(n),\qquad
w_X(n)=a_X(n)-b_X(n).
}
\tag{G0-source}
\]
The role of $b_X$ is to model the expected local density of the two affine prime forms;
$w_X$ is the discrepancy on which Type I/II cancellation is sought.

The elementary comparison algebra includes
\[
B(n)
=
\mathfrak S_2
\sum_{\substack{e\mid n\\e\ {\rm odd\ squarefree}}}
\prod_{\ell\mid e}\frac1{\ell-2},
\]
and for each fixed $d$,
\[
\lim_{Y\to\infty}
\frac1Y\sum_{m\le Y}B(dm)
=
\frac{d}{\phi(2d)}.
\tag{G0-Bmean}
\]
Thus one sign has mean $2d/\phi(2d)$ and the two-sign model has mean
$4d/\phi(2d)$ at the fixed-$d$ level.  Also
\[
B(n)\ll\log\log(3n),\qquad
|w_X(n)|\ll\log X,
\]
and, in the present source class,
\[
\sum_{n\asymp X}|w_X(n)|\tau(n)\ll X(\log X)^2.
\tag{G0-growth}
\]

\subsection{Gate-0 target package}

\begin{definition}[Gate-0 package]
Gate 0 is \emph{closed for the \#287 source at scale $X$} if the following seven items
hold with one common physical normalization.
\begin{enumerate}[label=\textup{G0.\arabic*},leftmargin=2.0em]
\item $a_X(n)\ge0$ and the comparison sequence $b_X$ belongs to the Ford--Maynard
ordinary sequence class used in the lower-bound theorem.
\item Prime mass:
\[
\sum_p b_X(p)\gg \frac{X}{\log X}.
\]
\item Growth/divisor control, including a bound of the form (G0-growth)
(or the exact stronger divisor-weight norm required by the chosen Ford invocation).
\item Comparison regularity: the physical $b_X$ satisfies the Ford--Maynard
condition $(b.2)$ uniformly over every coagulation region actually used downstream.
At the fixed-$d$ level this is compatible with (G0-Bmean); the required
uniform publication-level statement must be source-pinned.
\item Literal maximal/weighted Type I:
\[
\boxed{
\sum_{m\le X^{1/2-\varepsilon}}
\tau(m)^A
\max_I
\left|
\sum_{\substack{X/2<mn\le X\\n\in I}}
w_X(mn)
\right|
\ll_{A,W,\varepsilon}
X(\log X)^{-A}.}
\tag{G0-I}
\]
\item One fixed principal/small-conductor/exceptional-character expected-term
convention is used consistently by both the prime source and the comparison source.
\item Even, nonunit, prime-power and boundary strata created by the affine forms
$2n\pm1$ are either included in the source theorem or routed with a proved error.
\end{enumerate}
\end{definition}

\paragraph{Current verification.}
Items G0.1--G0.3 have elementary/source-level proofs in the dossier; the exact finite
problem statement and finite blocker are Lean-proved.  G0.4 and G0.5 are
\status{SOURCE-OPEN / EXTERNALLY-AUDITED TARGETS} at the exact publication normalization;
G0.6 is the same physical comparison convention that later appears in the factorial
balanced-seven and Pure5 comparison pins.

\paragraph{What Gate 0 closing buys.}
It supplies the non-parity hypotheses needed by the Ford lower-bound compiler and removes
the possibility that a later Type-II estimate is being applied to the wrong comparison
sequence.

\paragraph{What it does not buy.}
Gate 0 does not prove Type II, parity breaking, Gate~1A, Gate~1B, Gate~2, a large-$M$
witness, or Erd\H{o}s~\#287.

\section{Gate 1A: canonical common-weight Type-II provider}
\label{sec:gate1a-theory}

\subsection{Parameters, rows and physical source}

At an auxiliary scale $X_1$, set
\[
M=X_1^{1/3},\qquad
R=X_1^a,\qquad
L=X_1^b,
\]
\[
H=X_1^{a+2b-2/3},\qquad
K=\frac MR,\qquad
U=\frac LH,\qquad
D=\frac{L^2}{H}.
\tag{G1A-scales}
\]
Thus
\[
RK=M,\qquad DH=L^2,\qquad M^2H=RL^2.
\]
A row is
\[
e=(r,k,m),\qquad
r\asymp R\ {\rm prime},\quad
1\le|k|\le K,\quad
m\asymp M,\quad
m'=m+kr\asymp M.
\]
Choose the canonical affine forms $A_e(t),B_e(t)$ so that
\[
\boxed{m'A_e(t)-mB_e(t)=2k.}
\tag{G1A-det}
\]
For
\[
\rho_\ell(y)=\mathbf 1_{\ell\mid y}-\frac1\ell
\]
and bounded prime coefficients $b_p,d_q$, the \emph{one common physical-weight} source is
\[
\boxed{
C_e(b,d)
=
\sum_{t\in\mathbb Z}W_D(t)
\left(\sum_{p\asymp L}b_p\rho_p(A_e(t))\right)
\left(\sum_{q\asymp L}d_q\rho_q(B_e(t))\right).
}
\tag{G1A-source}
\]
The word ``common'' is literal: $W_D$ is one datum.  Arbitrary independently chosen
weights $W_{D,e}$ are a stronger source class and are not covered by the canonical theorem.

\subsection{Gate-1A target theorem}

\begin{theorem}[Canonical common-weight Gate-1A provider -- conditional form]
Assume the five source/analytic pins G1A-A through G1A-E below.  Then, uniformly for the
canonical all-$m$ row family $\mathcal E^\sharp$,
\[
\boxed{
\sum_{e\in\mathcal E^\sharp}|C_e(b,d)|^2
\ll
\frac{ML^4}{H}X_1^{o(1)}
\left(\sum_{p\asymp L}|b_p|^2\right)
\left(\sum_{q\asymp L}|d_q|^2\right).
}
\tag{G1A-target}
\]
The exact normalized variant used by the formal compiler is equivalent after the banked
$H^2$ normalization bridge.
\end{theorem}

The theorem is a \status{CONDITIONAL-COMPILER}: the finite/operator part is substantially
banked, but the physical source pins are not all inhabited.

\subsection{The five load-bearing pins}

\paragraph{G1A-A: Poisson--Bruhat/source transport.}
After the one-sided physical transform and reciprocity, every retained source state
$\xi$ must have moving-prime coefficient
\[
\boxed{c_r(\xi)=B_\xi\,F_\xi(r/R),}
\tag{G1A-PB}
\]
where $F_\xi$ is a uniformly smooth envelope generated from the one physical $W_D$,
not an arbitrary row weight.  The transport must preserve the physical/normalized energy
up to a square-summable source error.

\[
\status{STATUS: SOURCE-OPEN / CONDITIONAL PIN.}
\]
Without this pin, the later moving-prime participation argument would apply to a convenient
model rather than to the physical packet.

\paragraph{G1A-B: absolute source ceiling.}
The transported physical source must satisfy
\[
\boxed{
\mathcal T_{\rm abs}
\ll
M^2L^4X_1^{o(1)}.
}
\tag{G1A-abs}
\]
This is the absolute comparator against which the moving-family contraction is measured.

\[
\status{STATUS: SOURCE/ANALYTIC PIN.}
\]

\paragraph{G1A-C: fixed-state exclusion and prime participation.}
A fixed nonprojective state must not absorb the moving prime variable.  Smoothness of
$F_\xi$ produces a long plateau; the short-interval prime theorem then gives
\[
\boxed{N_\xi\gg R^{3/4-o(1)}}
\tag{G1A-part}
\]
participating prime rows for each retained state.  The finite BPP compiler converts this
participation into
\[
\boxed{
E_{\rm np}
\ll
R^{-1/2+o(1)}\mathcal T_{\rm abs}.
}
\tag{G1A-np}
\]
The project kernel-checks the finite participation-to-energy implication; it does not
formalize the entire physical smooth-envelope/short-interval input.

\noindent\textbf{Status.}
\status{FINITE COMPILER LEAN-PROVED}; the physical participation statement is
\status{EXTERNAL / SOURCE-OPEN}.

\paragraph{G1A-D: outer-root physical energy and projective routing.}
The abstract nonprojective energy must be identified with the physical Gate energy, and
the projective/axis/diagonal branch must be routed in the same source normalization.
The finite ratio-class algebra includes the state diagonal in the zero-curvature branch,
and the physical projective mass target is
\[
\boxed{
E_{\rm proj}
\ll
MHL^4X_1^{o(1)}.
}
\tag{G1A-proj}
\]
Exactly one outer square root converts (G1A-np) to an $R^{-1/4+o(1)}$ contraction.
At the three controlling vertices the signed margins are
\[
\boxed{\frac1{72},\qquad\frac1{24},\qquad\frac1{32}.}
\tag{G1A-margins}
\]

\noindent\textbf{Status.}
\status{FINITE SQUARE-ROOT / PROJECTIVE ALGEBRA BANKED}; the physical energy
identification is \status{SOURCE-OPEN}.

\paragraph{G1A-E: recombination/discarded-sector bound.}
After all source truncations, exceptional routes and recombinations, the physical error must
satisfy
\[
\boxed{
E_{\rm rec}
\ll
U^{-2+o(1)}M^2L^4,
}
\tag{G1A-rec}
\]
with the source-exact prefactor and without spending the same saving twice.  The exponent
margins at the controlling vertices are
\[
\boxed{\frac1{18},\qquad\frac1{18},\qquad\frac1{24}.}
\]

\noindent\textbf{Status.}
\status{EXPONENT BUDGET BANKED}; literal source recombination is \status{SOURCE-OPEN}.

\subsection{What the BPP mechanism proves and what it does not}

The finite BPP mechanism is a local-to-family compiler:
\[
\begin{gathered}
\text{smooth moving-prime participation}\\
\Downarrow\\
R^{-1/2+o(1)}\text{ family-energy contraction}\\
\Downarrow\\
R^{-1/4+o(1)}\text{ after one outer root}.
\end{gathered}
\]
It does \emph{not} prove the source transport, absolute source ceiling, physical energy
identification or recombination estimate.  Those are separate statements precisely so that
a reader can attack or falsify them independently.

\subsection{The source-glue/SB-$\nu$ problem}

The current source-glue question is whether the same physical packet can simultaneously
satisfy the common-weight representation, the BPP state representation and the final
routing representation with fibre multiplicity
\[
\boxed{N_\nu=X_1^{o(1)}}
\tag{G1A-SBnu}
\]
in the synchronized two-view model.  This is a source-identification frontier rather than
a new exponent computation.

Finally,
\[
\boxed{\texttt{GATE1A\_REQUIRED\_FOR\_287 = UNKNOWN}.}
\]
The Ford-generated packet census, not programme nomenclature, decides whether a needed
\#287 packet truly belongs to the canonical common-weight source class.

\section{Gate 1B: signed/M\"obius parity-breaking provider}
\label{sec:gate1b-theory}

\subsection{Native centering, sign and complementary-divisor geometry}

For coprime $d,p$,
\[
\boxed{
\rho_{dp}(y)
=
\rho_d(y)\rho_p(y)
+\frac{\rho_d(y)}p
+\frac{\rho_p(y)}d.
}
\tag{G1B-center}
\]
If $q$ is squarefree and $p\mid q$ is the distinguished prime, then
\[
\boxed{\mu(q/p)=-\mu(q).}
\tag{G1B-sign}
\]
This is the sign resource that must not be squared away accidentally.

The native complementary-divisor source exposes
\[
\boxed{q\ell-uv=2.}
\tag{G1B-det}
\]
This determinant-two geometry is separate from the Gate-1A determinant $2k$ and from the
\#287 $\mu*\log$ determinant-one line.

\subsection{The banked one-completion source}

On dyadic windows $d\sim D$ and $p\sim P$, the authoritative one-completion coefficient is
\[
\boxed{
\beta_{D,P}(d,p)=\mu_D(d)\Lambda_P(p),
}
\tag{G1B-beta}
\]
where $\mu_D,\Lambda_P$ are the physical truncated M\"obius and prime weights on the stated
windows.  The public formal challenge proves the separate linearity/bilinearity of this
coefficient and an exact level-of-distribution budget requiring
\[
\boxed{
\theta\ge\frac12+\frac1{104}.
}
\tag{G1B-104}
\]
It also proves the implication from an inhabited source-to-convolution-BV interface to the
abstract Gate conclusion.  Producing that interface from the literal physical source remains
open.  This is a valid formal provider socket, but it is no longer the first research-level
analytic obstruction.

\subsection{The Motohashi line}

Fix $u,\ell$ and a base solution to
\[
\boxed{uv_0+2=\ell z_0.}
\tag{G1B-base}
\]
All solutions on the complementary-divisor line are
\[
\boxed{
v_t=v_0+\ell t,\qquad
z_t=z_0+ut.
}
\tag{G1B-line}
\]
Let $H$ be the physical line length and write $e_H(x)=e^{2\pi ix/H}$.
For the exact one-completion coefficient and a smooth cutoff $W_1$ define
\[
\boxed{
\mathcal Z_{u,\ell}(k)
=
\sum_t
\beta_{D,P}(z_t)\,
W_1(t/H)e_H(-kt).
}
\tag{G1B-Z}
\]
For the five centred defect factors
\[
b_5=\delta_1*\delta_2*\delta_3*\delta_4*\delta_5
\]
and a second physical cutoff $W_2$, define
\[
\boxed{
\mathcal B_{u,\ell}(k)
=
\sum_t
b_5(v_t)\,
W_2(t/H)e_H(-kt).
}
\tag{G1B-B}
\]
The Motohashi pairing is
\[
\boxed{
\mathcal M_{u,\ell}
=
\frac1H
\sum_k
\mathcal Z_{u,\ell}(k)\,
\overline{\mathcal B_{u,\ell}(k)}.
}
\tag{G1B-M}
\]
The current supplied research ledger reports the exact A/B/C source pin, the twisted-defect
A/B/C reduction and the fivefold Motohashi iteration as passed source algebra; the
fixed-polylogarithmic-$k$ interior is closed modulo the comparison convention.

\subsection{Endpoint residue coefficient and diagonal/off-diagonal split}

At the endpoint define, for a unit residue $a\bmod\ell$,
\[
\boxed{
A_{\ell,a,k}
=
\sum_{\substack{u\equiv-2a^{-1}\pmod\ell}}
a_4(u)\,\mathcal Z_{u,\ell}(k),
}
\tag{G1B-A}
\]
where $a_4$ is the four-defect coefficient inherited from the physical source.
The endpoint Cauchy energy is
\[
\boxed{
\mathcal E_{\rm end}(k)
=
\sum_{\ell\sim L}
\sum_{a\in(\mathbb Z/\ell\mathbb Z)^\times}
|A_{\ell,a,k}|^2.
}
\tag{G1B-Eend}
\]
Opening the square gives
\[
\mathcal E_{\rm end}(k)
=
\sum_{\ell,a}
\sum_{\substack{u_1,u_2\equiv-2a^{-1}\,(\ell)}}
a_4(u_1)\overline{a_4(u_2)}
\mathcal Z_{u_1,\ell}(k)
\overline{\mathcal Z_{u_2,\ell}(k)}.
\]
The diagonal is $u_1=u_2$.  The off-diagonal is parametrized by
\[
\boxed{u_2=u_1+j\ell,\qquad j\ne0.}
\tag{G1B-j}
\]
For corresponding line parameters $t_1,t_2$,
\[
\boxed{
z_2(t_2)-z_1(t_1)
=
u(t_2-t_1)+j\,v_{t_2},
}
\tag{G1B-off-source}
\]
and, after opening $\beta=\mu_D\Lambda_P$,
\[
\boxed{
d_2p_2-d_1p_1
=
u(t_2-t_1)+j\,v_{t_2}.
}
\tag{G1B-off-beta}
\]
These identities are the literal arithmetic object behind the current endpoint
off-diagonal problem.

The latest supplied research ledger records
\[
\boxed{
\text{endpoint }u\text{-diagonal}
=
\text{power closed with }X^{-1/12+o(1)}.
}
\tag{G1B-diag-status}
\]

\subsection{Current first Gate-1B analytic theorem}

Let $\mathfrak T_{1B}(X)$ denote the physical target scale of the Gate-1B parent after
restoring the dyadic prefactor from the source decomposition.  The current source materials
provided to this editorial run specify the off-diagonal object (G1B-Eend)--(G1B-off-beta)
but do not print the final numerical parent prefactor in a form that can be independently
reconstructed here.  We therefore do \emph{not} invent it.

The current theorem slot is:
\[
\boxed{
\sum_{\text{endpoint dyadic packets}}
\mathcal E_{\rm end}^{\rm off}
=
o\!\left(\mathfrak T_{1B}(X)\right),
}
\tag{G1B-off-target}
\]
with the exact physical prefactor/normalization to be copied verbatim from the controlling
Gate-1B source before publication.  Programme label:
\[
\boxed{\begin{gathered}\text{\ttfamily RANKONE-ENDPOINT-}\\\text{\ttfamily U-OFFDIAG45}\end{gathered}}
\]
Its source formula is exact; the missing printed normalization in this candidate is
\status{SOURCE-BLOCKED NORMALIZATION}, and the cancellation itself is
\status{ANALYTIC-OPEN}.  This explicit split is preferable to guessing a power.

\subsection{Downstream Gate-1B theorem objects}

The first endpoint theorem is not the whole Gate.  The current dependency chain is:

\paragraph{PURE5 comparison/main-term pin.}
The five-defect physical source and its comparison source must use the same principal,
small-conductor and exceptional-character convention.
\[
\status{\path{PURE5-COMPARISON-MAINTERM-PIN}: SOURCE-OPEN.}
\]
Closing an oscillatory estimate with the wrong main term does not close the physical packet.

\paragraph{High-$k$ sector.}
The same Motohashi pairing (G1B-M) restricted to the complementary high-frequency
range defines \path{RANKONE-HIGHK45}.  Its target is the Gate-1B parent bound on that
frequency sector.
\[
\status{ANALYTIC-OPEN.}
\]

\paragraph{Pure-five signed distinguished-prime packet.}
After recombining interior, endpoint and high-$k$ sectors, the source is the signed
five-centred-defect distinguished-prime packet, schematically
\[
\sum \mu_D(d)\Lambda_P(p)
\prod_{i=1}^5\delta_i(\cdot)
\times(\text{physical cutoffs}),
\]
with all congruence and gcd restrictions inherited from (G1B-det).
The literal final prefactor is part of the source packet.
\[
\status{\path{PURE5-DP-SIGNED45}: DOWNSTREAM ANALYTIC/REASSEMBLY OPEN.}
\]

\paragraph{Lower defect orders.}
The $|J|=4,3,2,1$ packets are separate centred-defect strata obtained by selecting which
defect factors remain nontrivial.  They do not follow automatically from the Pure5 theorem;
each requires its own source routing or a proved monotonicity/reduction theorem.

\paragraph{Near-primitive and conductor complement.}
The near-primitive $r=1$ full-conductor slice is denoted
\texttt{NEARPRIM-DP-SIGNED45}.  The $r>1$ family, the
$(C_\ast,C_{\rm nw})$ transition, and proper-divisor recursion must then restore the
remaining conductor ranges without losing the M\"obius sign.

\paragraph{QK56 parent.}
The full covariance parent has the schematic but literal two-modulus form
\[
\boxed{
\mathcal C_{\rm QK}
=
\sum_{q_1,q_2}
\sum_{h_1,k_1,h_2,k_2}
A(q_1,q_2;h_1,k_1,h_2,k_2)
Q(q_1;h_1,k_1)
\overline{Q(q_2;h_2,k_2)}.
}
\tag{QK56-parent}
\]
The branch $q_1=q_2$ is the same-modulus character Gram; $q_1\ne q_2$ is the
cross-modulus/CRT branch.  The canonical zero mode and substantial finite/source
decompositions are banked, but \path{QK56-EXHAUSTIVENESS} must prove that all physical
children are covered.  QK56 is downstream of the rank-one endpoint in the current
closure order, not the present first analytic open.

\paragraph{Shifted $TT^\ast$ and finite reassembly.}
The final shifted source must be recovered before any support enlargement; all zero/nonunit,
collision and repeated-prime strata must be routed; then the finite source compiler
reassembles the leaves into the Gate-1B parent.

Only after these steps may one write \texttt{GATE1B CLOSED}.

\section{Ford--Maynard generated packets and proof-specific Type II}
\label{sec:ford-theory}

\subsection{Published universal wrapper}

The Ford--Maynard lower-bound framework uses a universal Type-II hypothesis.  In the
present $1/6$ specialization a safe sufficient form is
\[
\boxed{
\left|
\sum_{\substack{(X/2)^{\varepsilon}<m\le X^{1/6-\varepsilon}\\X/2<mn\le X}}
\xi_m\kappa_n w_X(mn)
\right|
\ll_A X(\log X)^{-A}
}
\tag{FM-universal-II}
\]
for every allowed divisor-bounded coefficient pair.  The project does \emph{not} claim
that Ford--Maynard publish the weaker proof-specific theorem described next.

\subsection{What their proof actually generates}

The Ford arithmetic source has
\[
G(m;n)
=
\mathbf 1_{m\le n^\gamma}\lambda(m_1)\,
g(\mathbf v(m_2;n)),
\qquad
\boxed{H(n)=\sum_{d\mid n}G(d;n),}
\tag{FM-GH}
\]
with the canonical rough/smooth split $m=m_1m_2$.
After finite decomposition and condition removal, Proposition~7.22 reaches a fixed-depth
product sum (equation~(7.23) in the source)
\[
\sum_{n=n_1\cdots n_N\asymp x}
w_n\,x_1(n_1)\cdots x_N(n_N),
\tag{FM-7.23}
\]
and chooses a product $\prod_{e\in E}n_e$ in the Type-II range.
Grouping the selected factors and their complement gives
\[
\boxed{
\sum_{n=n'n''\asymp x}
Y_1(n')Y_2(n'')w_{n'n''}.
}
\tag{FM-generated-II}
\]
The published proof then invokes the universal Type-II hypothesis.

Hence an invocation-by-invocation replacement theorem would be enough for this proof:
\[
\boxed{
\begin{gathered}
\text{every actual generated (FM-generated-II) packet}\\
\text{has a proved source-exact provider}\\
\Downarrow\\
\text{every Type-II invocation in Proposition 7.22 is covered}
\end{gathered}}
\tag{FM-replacement-logic}
\]
This implication additionally requires a source-exhaustive census and a replacement
certificate; it is not obtained merely by identifying some packets.

\subsection{The three source obligations}

\begin{definition}[SOURCE-ADAPTER45]
For every generated packet, print the literal operator, physical ranges, coefficient
class, weight dependence, centering/expected term, gcd/collision restrictions, original
prefactor and normalized target, and map it to a proved provider without changing the
source class.
\end{definition}

\begin{definition}[FM-TO-GATE-COORDINATE-CENSUS45]
Enumerate every Type-II packet actually generated by the Ford decomposition, with enough
coordinate information to decide whether it belongs to Gate~1A, Gate~1B or a direct
exceptional/standard class.
\end{definition}

\begin{definition}[$C_{\rm FM}$ replacement certificate]
Verify invocation by invocation that the generated census covers every use of the universal
Type-II hypothesis in the chosen Ford proof, with no omitted packet, no multiplicity loss
and no double-spending of an analytic saving.
\end{definition}

All three remain open source/reassembly obligations.

\subsection{Provider-facing generated families}

\begin{longtable}{p{0.19\textwidth}p{0.20\textwidth}p{0.20\textwidth}p{0.30\textwidth}}
\toprule
Family & Literal source/scale & Candidate provider & Why it matches / what is missing\\
\midrule
\endfirsthead
\toprule
Family & Literal source/scale & Candidate provider & Why it matches / what is missing\\
\midrule
\endhead
M\"obius singleton &
one selected $m<X^{1/6}$ factor carries M\"obius weight; complement remains generated &
Gate~1B signed source &
M\"obius sign is in the correct physical slot only after a literal completion dictionary;
generated outer-factor stability is still required.\\
Model singleton &
selected factor is model/Mellin-like; M\"obius carrier may remain in the complement &
QK56 or direct affine completion &
provider cannot be inferred from the word ``singleton''; native packet dictionary is open.\\
Balanced-seven V-cell &
seven fixed prime boxes at $Y=X^{1/7}$ &
factorial-$\Omega_7$ source, then the fixed-degree-seven attack &
factorial Euler encoding is exact; endpoint/3221 analytic theorem and comparison remain.\\
Pure5/high-order family &
five centred defect factors with distinguished-prime/M\"obius source &
Gate~1B Motohashi/rank-one chain &
current first analytic object is the endpoint off-diagonal; high-$k$ and downstream
defect/conductor families remain.\\
Balanced $R_9$ &
nine-coordinate Ford-generated leakage cell &
fixed-certificate/FCL lane &
finite binomial $70$ is proved; literal physical $G(d;n)$ specialization is source-open.\\
Common-weight packets &
actual generated source must preserve one physical $W_D$ after scale conversion &
Gate~1A &
requires exact scale/source dictionary and all five Gate-1A pins.\\
Remaining fixed-$g_\ast$ packets &
not yet exhaustively censused &
TBD from literal source &
the finite census/provider assignment itself is still open.\\
\bottomrule
\end{longtable}

\section{Balanced $R_9$, fixed-certificate leakage and FCL}
\label{sec:r9-fcl-theory}

The balanced order-nine finite identity is
\[
\boxed{
\sum_{j=0}^{4}(-1)^j\binom9j=70.
}
\tag{R9-70}
\]
This is \status{PROVED-FINITE}; by itself it is not the physical Ford statement
$H(n)=70$.

To obtain $H(n)=70$ one must prove that the literal Ford specialization of $G(d;n)$ on the
balanced $R_9$ cell selects exactly the corresponding low/high divisor orders and has the
required value of $g(\varnothing)$.  That is
\[
\boxed{\text{literal }G(d;n)\text{ specialization: SOURCE-OPEN}.}
\tag{R9-G-open}
\]
Therefore balanced $R_9$ is a \emph{Ford-generated leakage candidate}, not a completed
leakage theorem.

For one fixed positive Ford certificate $g_\ast$, define its arithmetic weight
\[
H_\ast(n)=\sum_{d\mid n}G_\ast(d;n).
\]
Let $\mathcal P$ be the desired prime region, $\mathcal N_1$ the certificate-sign composite
region, $\mathcal N_2$ the exceptional collar, and
\[
\mathcal U_\ast
=
(X/2,X]\setminus(\mathcal P\cup\mathcal N_1\cup\mathcal N_2).
\]
The fixed-certificate leakage is
\[
\boxed{
\mathcal L_\ast(X)
=
\sum_{n\in\mathcal U_\ast}w_X(n)H_\ast(n).
}
\tag{FCL}
\]
The $N_2$ error is
\[
\mathcal E_{N_2}(X)
=
\sum_{n\in\mathcal N_2}w_X(n)H_\ast(n).
\]
If the total correlation, leakage, $N_2$ error and comparison error are all small relative
to the positive certificate margin $C_\ast(\varepsilon_\ast)$, then the fixed-certificate
transference gives positive affine prime mass.  A sufficient quantitative form is
\[
|\mathcal L_\ast(X)|+|\mathcal E_{N_2}(X)|
\le
\frac{C_\ast(\varepsilon_\ast)}4\,B_X,
\tag{FCL-margin}
\]
together with the corresponding total-correlation/comparison bounds.
This is a conditional compiler; it does not itself prove FCL.

Balanced $R_9$ may contribute to $\mathcal L_\ast$ once the source obligation (R9-G-open) is discharged,
but closing that one cell does not close the full leakage family.

\section{Gate 2: source-specific Ford--Maynard lower-bound compiler}
\label{sec:gate2-theory}

Fix one compatible
\[
P_{\varepsilon_\ast}
=
\left(
\frac12-\varepsilon_\ast,\,
\varepsilon_\ast,\,
\frac16-2\varepsilon_\ast
\right).
\]
Gate~2 is the theorem-shaped endgame package:

\begin{theorem}[Gate-2 lower-bound compiler -- conditional]
Assume:
\begin{enumerate}[label=\textup{G2.\arabic*},leftmargin=2.0em]
\item the Gate-0 source/comparison/Type-I package;
\item the literal Ford $G(d;n)$/$H(n)$ source and its sign properties;
\item a source-specific aggregate $N_2$ estimate;
\item one common fixed $\varepsilon_\ast>0$ compatible with every source range;
\item either the universal Type-II theorem (FM-universal-II), or a complete
generated-packet provider package together with
\path{SOURCE-ADAPTER45},
\path{FM-TO-GATE-COORDINATE-CENSUS45} and $C_{\rm FM}$;
\item the ordinary/bounded-class firewall and every comparison/exceptional-character
normalization needed by the chosen Ford theorem;
\item effective versions of all inputs at a computable threshold.
\end{enumerate}
Then the Ford--Maynard lower-bound geometry yields positive affine mass of the form
\[
\boxed{
\sum_{p\in\mathcal P}a_X(p)
\ge
\bigl(C_\ast(\varepsilon_\ast)+o(1)\bigr)B_X
\gg \frac{X}{\log X},
}
\tag{G2-positive}
\]
for the physical \#287 source and the chosen certificate.
\end{theorem}

\[
\status{STATUS: CONDITIONAL-COMPILER / NOT ACTIVATED.}
\]

The theorem still does not by itself produce \texttt{WindowPairSupply}.  One additional effective witness
extraction theorem must convert the positive affine mass into a Sophie/safe-prime or
prime-power blocker configuration, or directly into the $\WPS(M)$ predicate of the finite
compiler.  That final conversion/effectivity step remains open.

\section{Current \#287 specialization: $\mu*\log$, factorial $\Omega_7$, and the 3221 attack}
\label{sec:287-specialization-theory}

The source-minimal identity used in the current \#287 branch is
\[
\boxed{
\Lambda(2mn+s)
=
\sum_{qr=2mn+s}\mu(q)\log r,
\qquad s=\pm1.
}
\tag{287-mulog}
\]
The latest supplied V15 replay reports the exact three-way split into small-$q$,
small-$r$, and hard pairs as kernel-checked, together with the determinant-one line
\[
\boxed{qr-2mn=s.}
\]
For fixed $r,m$ and one solution,
\[
q=q_0+2mt,\qquad n=n_0+rt.
\]

\subsection{Factorial Euler seven-box source}

For the balanced-seven cell
\[
Q=X^{3/5},\qquad Y=X^{1/7},
\]
let
\[
a_{\mathbf z}(p)=\frac17\sum_{i=1}^7z_i\omega_i(p),
\qquad
\boxed{
F_{\mathbf z}(p^e)=\frac{a_{\mathbf z}(p)^e}{e!}.
}
\tag{Omega7-factorial}
\]
Then the degree-$(1,\ldots,1)$ torus coefficient exactly recovers the ordered seven-labelled
prime source, including repeated primes:
\[
\boxed{
7^7[z_1\cdots z_7]F_{\mathbf z}(n)
=
\sum_{\substack{p_1\cdots p_7=n\\p_i\ {\rm ordered}}}
\prod_{i=1}^7\omega_i(p_i).
}
\tag{Omega7-coeff}
\]
The latest supplied V16 factorial replay reports the general/seven-slot polarization,
repeated-prime fibre cardinality, formal local Euler algebra and expected-term linearity as
Lean-proved in an $8123$-job build.

The exact balanced-seven source therefore reduces to
\[
\int_{\mathbb T^7}\overline{z_1\cdots z_7}
\sum_{q\sim X^{3/5}}\mu(q)
\Delta_{\mathcal E}(F_{\mathbf z};X,q,-s\overline2)\,d\mathbf z
\]
plus the physical comparison identification.  The exact broad endpoint theorem is
\[
\boxed{\begin{gathered}\text{\ttfamily AFFINE287-FACTORIAL-OMEGA7-}\\\text{\ttfamily SIGNED-ENDPOINT45}\end{gathered}}
\tag{Omega7-endpoint}
\]
It is strictly narrower than a BV theorem for all class-$\mathcal C$ functions, an
absolute modulus average, a theorem uniform in every $\mathbf z$, or a general
smooth-number theorem.

\subsection{Why the direct existing Pascadi proof route does not splice}

In the audited current v2 proof architecture, a small parameter $\eta$ is arranged so that
\[
Q\le x^{5/8-100\eta},
\qquad
y\le x^\eta.
\]
At $Q=x^{3/5}$,
\[
\eta\le\frac1{4000},
\]
whereas $y=x^{1/7}$ requires
\[
\eta\ge\frac17.
\]
Hence the direct written proof route is nonclosing at the required smoothness scale.  This is
a parameter/dictionary audit, not a theorem ruling out future extensions.

\subsection{The current 1+2+2+2 / 3221 research reduction}

The latest supplied Pro reduction keeps the source at fixed degree seven and groups the
seven prime boxes as
\[
\boxed{1+2+2+2.}
\]
Use the scale ledger
\[
\boxed{
E=X^{1/7},\qquad
M=N=L=X^{2/7},\qquad
Q=X^{3/5}.
}
\tag{3221-scales}
\]
The ordinary three-block comparison sees
\[
\frac47-\frac35=-\frac1{35},
\]
so the naive $Q<NL$ geometry misses by $X^{1/35}$.
Retaining the actual seventh box inside Cauchy changes the controlling comparison to
\[
\boxed{Q<ENL}
\]
with exponent margin
\[
\boxed{\frac57-\frac35=\frac4{35}.}
\tag{3221-margin}
\]
The source-assisted diagonal audit gives an amplitude saving
\[
\boxed{X^{-2/35+o(1)}.}
\tag{3221-saving}
\]
After opening the two copies, the literal difference relation is
\[
\boxed{
e_1n_1\ell_1-e_2n_2\ell_2=qt,
\qquad
|t|\ll X^{4/35+o(1)}.
}
\tag{3221-diff}
\]
The line-length ledger is
\[
\boxed{
H=X^{11/35},
\qquad
EH=X^{16/35}<Q.
}
\tag{3221-H}
\]

The current proposed narrower source socket is
\[
\boxed{\begin{gathered}\text{\ttfamily 3221-DI-KUZNETSOV-}\\\text{\ttfamily LITERAL-SPLICE45}\end{gathered}}
\tag{3221-label}
\]
Its theorem is not ``prove Kuznetsov''.  It is the source-dictionary/analytic statement:
after the above post-dispersion decomposition, identify the literal off-diagonal
(3221-diff) with a DI/Kuznetsov-type bilinear form \emph{without changing the physical
weights, M\"obius sign, prefactor or multiplicity}, and obtain a fixed-power saving sufficient
to recover at least the source-assisted margin (3221-saving).  If such a literal splice
is proved, it is a narrower sufficient route toward (Omega7-endpoint).

\[
\status{STATUS: RESEARCH CANDIDATE / SOURCE-DICTIONARY + ANALYTIC OPEN.}
\]
The supplied materials do not contain a completed hostile audit of this splice, so it is not
promoted above the exact factorial endpoint theorem.

\section{Verification manual: what is actually proved}
\label{sec:verification-manual}

\subsection{Table 1: machine-checked objects}

{\footnotesize\begin{longtable}{p{0.20\textwidth}p{0.22\textwidth}p{0.23\textwidth}p{0.20\textwidth}}
\toprule
Object / declaration & File or replay & Exact content & Does not prove\\
\midrule
\endfirsthead
\toprule
Object / declaration & File or replay & Exact content & Does not prove\\
\midrule
\endhead
\parbox[t]{0.19\textwidth}{\raggedright\scriptsize\ttfamily no\_Erdos287\\Counterexample\_\\of\_max\_le\_4e9} &
\path{FiniteRangeExtension.lean} &
No exact counterexample with maximum denominator at most $4\cdot10^9$. &
No large-$M$ theorem.\\
\parbox[t]{0.19\textwidth}{\raggedright\scriptsize\ttfamily WindowPairSupply;\\no\_Erdos287\\Counterexample\_\\of\_closure} &
\path{ClosureInputs.lean} / end-to-end status bank &
Exact conditional compiler from effective eventual supply plus finite range. &
No inhabitant of eventual supply.\\
\path{vaughan_identity_exact}; source compilers &
V14 affine Vaughan files &
Exact alternative source decomposition and structural algebra. &
No Type-II estimate.\\
\parbox[t]{0.19\textwidth}{\raggedright\scriptsize\ttfamily
vonMangoldt\_eq\_\\mobius\_mul\_log;\\muLog\_affine\_\\pointwise} &
latest supplied V15 replay:
\path{AffineMuLogIdentity.lean} &
$\Lambda=\mu*\log$ and its affine specialization. &
No cancellation.\\
three-way pair partition / \texttt{muLogHard} &
\path{AffineMuLogHardSource.lean} &
Exact disjoint small-$q$/small-$r$/hard decomposition. &
No smallness of any piece.\\
determinant-one line family &
latest supplied V15 replay:
\path{AffineMuLogLine.lean} &
Coprimality, affine parametrization and multiplicity-one uniqueness. &
No dispersion estimate.\\
balanced-seven $-20$ &
\path{BalancedSevenFinite.lean} in latest supplied V15 replay &
Finite alternating-binomial coefficient. &
No physical Ford source specialization.\\
squarefree labelled polarization &
\parbox[t]{0.21\textwidth}{\raggedright\scriptsize latest supplied V15 replay:\\{\ttfamily BalancedSeven\\Polarization.lean}} &
Exact squarefree coefficient extraction and normalization. &
Superseded as controlling by factorial encoding.\\
\parbox[t]{0.19\textwidth}{\raggedright\scriptsize\ttfamily OMEGA7-FACTORIAL-\\EULER-\\POLARIZATION45} &
latest supplied V16 factorial replay; exact public filename synchronization pending &
General/seven-slot factorial polarization, including repeated-prime fibre accounting. &
No endpoint distribution estimate.\\
\parbox[t]{0.19\textwidth}{\raggedright\scriptsize\ttfamily POLARIZED-\\EXPECTED-TERM-\\LINEARITY45} &
latest supplied V16 factorial replay &
Coefficient extraction commutes with every linear expected-term operator. &
No physical identification $\mathfrak M^{\rm fac}=\mathfrak M^{\rm phys}$.\\
Pascadi $(3/5,1/7)$ parameter ledger &
latest supplied V16 factorial replay &
Formal rational inequality $\eta\le1/4000$ versus $\eta\ge1/7$. &
No theorem against future analytic extensions.\\
\bottomrule
\end{longtable}

The supplied V16 factorial report states \texttt{lake build}: $8123$ jobs, zero errors; no
Lean-code occurrences of \texttt{sorry}, \texttt{admit}, user \texttt{axiom},
\texttt{opaque}, \texttt{unsafe}, \texttt{native\_decide} or
\texttt{@[implemented\_by]}; principal axiom reports use only the standard logical
principles \texttt{propext}, \texttt{Classical.choice} and \texttt{Quot.sound}.
These are replay-status claims supplied with the project; public aggregate synchronization
is tracked separately.

\subsection{Table 2: externally audited / mathematical inputs}

\begin{longtable}{p{0.30\textwidth}p{0.22\textwidth}p{0.38\textwidth}}
\toprule
Input & Status & Role / limitation\\
\midrule
\endfirsthead
\toprule
Input & Status & Role / limitation\\
\midrule
Huxley short-interval prime theorem &
\status{EXTERNALLY-AUDITED} &
Supplies prime participation once the Gate-1A physical smooth envelope is source-pinned.\\
Ford--Maynard lower-bound framework &
\status{EXTERNALLY-AUDITED} &
Publishes the universal Type-II sufficient wrapper and lower-bound geometry; does not
publish the project's proof-specific replacement theorem.\\
Drappeau--Granville--Shao class-$\mathcal C$ distribution &
\status{EXTERNALLY-AUDITED / RANGE MISMATCH} &
Relevant multiplicative-function technology; does not literally give the exact
$Q=X^{3/5}$ endpoint at the required source.\\
Pascadi smooth-number proof architecture &
\status{EXTERNALLY-AUDITED / DIRECT ROUTE NONCLOSING} &
At $Q=X^{3/5}$ its written parameter architecture forces $\eta\le1/4000$, incompatible
with $y=X^{1/7}$.\\
Latest Motohashi/rank-one Gate-1B reductions &
\status{RESEARCH-AUDITED, NOT PUBLIC THEOREMS} &
Source reductions supplied after the V15 snapshot; first analytic object is the endpoint
off-diagonal, with exact parent normalization still to be copied into the publication draft.\\
Latest 3221 reduction &
\status{RESEARCH CANDIDATE} &
Narrows the factorial endpoint to a post-dispersion fixed-degree-seven source socket if
the DI/Kuznetsov literal splice survives hostile audit.\\
\bottomrule
\end{longtable}

\subsection{Table 3: open source identifications}

{\footnotesize\begin{longtable}{p{0.30\textwidth}p{0.24\textwidth}p{0.31\textwidth}}
\toprule
Source object & Status & What must be printed/proved\\
\midrule
\endfirsthead
\toprule
Source object & Status & What must be printed/proved\\
\midrule
Gate-0 comparison $(b.2)$/Type-I normalization &
\status{SOURCE-OPEN} &
Uniform physical progression/coagulation convention and literal weighted/maximal theorem.\\
Gate-1A PB/source glue and physical energy pin &
\status{SOURCE-OPEN} &
Exact transport, common-weight preservation, absolute ceiling, participation state, energy
identification, projective routing and recombination.\\
Gate-1B endpoint physical target prefactor &
\status{SOURCE-BLOCKED NORMALIZATION} &
Copy the literal parent prefactor/target accompanying (G1B-off-target); do not infer it
from the programme label.\\
\path{PURE5-COMPARISON-MAINTERM-PIN} &
\status{SOURCE-OPEN} &
Principal/small-conductor/exceptional-character physical main-term convention.\\
Ford generated-(7.23) census &
\status{SOURCE-OPEN} &
Every actual packet, scale, coefficient class, weight law, multiplicity and provider.\\
\path{SOURCE-ADAPTER45} / $C_{\rm FM}$ &
\status{SOURCE-OPEN} &
Literal Gate/provider dictionary and invocation-by-invocation replacement certificate.\\
Balanced-$R_9$ physical $G(d;n)$ specialization &
\status{SOURCE-OPEN} &
Show that the physical Ford cell really realizes the divisor-order split leading to $70$.\\
Factorial/physical comparison &
\status{SOURCE-OPEN} &
Prove $\mathfrak M^{\rm fac}_{7,s}=\mathfrak M^{\rm phys}_{7,s}$ in the exact local-density
and exceptional-character convention.\\
3221 DI/Kuznetsov completed-source dictionary &
\status{SOURCE-OPEN} &
Identify the post-dispersion source with the intended theorem socket with no loss of sign,
weights, prefactor or multiplicity.\\
\bottomrule
\end{longtable}

\subsection{Table 4: open analytic theorems in current dependency order}

{\footnotesize\begin{longtable}{p{0.22\textwidth}p{0.31\textwidth}p{0.32\textwidth}}
\toprule
Open theorem & Exact task & What remains even if proved\\
\midrule
\endfirsthead
\toprule
Open theorem & Exact task & What remains even if proved\\
\midrule
\path{3221-DI-KUZNETSOV-LITERAL-SPLICE45} &
Prove the literal post-dispersion 1+2+2+2 source admits a DI/Kuznetsov-type estimate with a
fixed-power saving sufficient for the source-assisted margin (3221-saving). &
Physical comparison and the rest of the factorial endpoint/source census; status remains a
research candidate until the source dictionary is hostile-audited.\\
\path{AFFINE287-FACTORIAL-OMEGA7-SIGNED-ENDPOINT45} &
Bound the signed degree-$(1,\ldots,1)$ torus coefficient of the $F_{\mathbf z}$ discrepancy
by $o(X/\log X)$. &
Physical comparison match; other V-cells and M-branch.\\
\path{RANKONE-ENDPOINT-U-OFFDIAG45} &
Bound the off-diagonal endpoint energy (G1B-Eend)--(G1B-off-beta) relative to
the physical Gate-1B parent target. &
Pure5 comparison, high-$k$, lower defects, conductor recursion, QK56 and reassembly.\\
\path{RANKONE-HIGHK45} &
Control the high-frequency part of (G1B-M) in the physical source normalization. &
Pure5/lower-defect/conductor/QK56 chain.\\
Remaining generated packet providers &
Prove each uncensused Ford-generated packet has a valid Gate~1A, Gate~1B or direct provider. &
Replacement certificate, FCL, Gate~2.\\
FCL &
Control the leakage sum (FCL) and the associated total-correlation/$N_2$/comparison errors within the
fixed certificate margin. &
Gate~2 source package/effectivity and large-$M$ witness extraction.\\
$N_2$/Gate-2 analytic inputs &
Prove the source-specific aggregate exceptional-collar estimate and all remaining
lower-bound hypotheses at one $\varepsilon_\ast$. &
Positive mass still must be converted effectively to \texttt{WindowPairSupply}.\\
eventual \texttt{WindowPairSupply} &
Prove $\WPS(M)$ for every sufficiently large $M$, or prove a stronger effective witness
implying it. &
Nothing analytic: the existing finite Lean compiler then gives Erd\H{o}s~\#287.\\
\bottomrule
\end{longtable}

\section{Current numbered route to Erd\H{o}s \#287}
\label{sec:numbered-route}

The current literal route, subject to future simplification, is:
\begin{enumerate}[leftmargin=2.1em]
\item finish or supersede the current fixed-degree-seven/3221 socket;
\item prove the factorial physical comparison identification;
\item complete the remaining balanced-seven/V-cell census;
\item control the M\"obius singleton and model singleton packets;
\item control the $k=0$ smooth-parity family;
\item complete the remaining fixed-$g_\ast$ packet census;
\item assign every generated packet to a proved Gate~1A, Gate~1B or direct provider;
\item complete the generated-(7.23) replacement certificate $C_{\rm FM}$;
\item close the full fixed-certificate leakage/correlation package;
\item prove the source-specific $N_2$ estimate;
\item activate Gate~2 at one compatible $\varepsilon_\ast$ and obtain positive affine mass;
\item make the lower-bound/witness extraction effective;
\item prove eventual $\WPS(M)$ (or a stronger effective witness);
\item invoke the existing finite Lean compiler.
\end{enumerate}
A future theorem may collapse several steps; the list is a current dependency audit, not a
claim of logical minimality.

\section{How an external reviewer can falsify or improve the programme}
\label{sec:review-invitation}

A serious review is most useful if it attacks one of the following:
\begin{enumerate}[leftmargin=2.0em]
\item a source dictionary that changes the physical packet;
\item a theorem-slot mismatch between an actual generated packet and a claimed provider;
\item a missing normalization or prefactor;
\item a multiplicity or support-enlargement loss;
\item double-spending of a M\"obius sign or analytic saving;
\item an incorrect comparison/principal/exceptional-character main term;
\item an incomplete Ford-generated packet census;
\item an invalid invocation-by-invocation replacement of Ford--Maynard's universal Type II;
\item an endpoint/effectivity gap;
\item a mismatch between a claimed Lean status and the actual kernel/public bank.
\end{enumerate}
Failure of any conditional provider does \emph{not} invalidate the machine-checked finite
theorem through $4\times10^9$; it invalidates only the corresponding large-$M$ route.

\section{Glossary of recurring programme terms}
\label{sec:glossary}

\begin{description}[leftmargin=3.4cm,style=nextline]
\item[Gate 0] The physical \#287 source/comparison/Type-I package.
\item[Gate 1A] Canonical common-weight moving-source Type-II provider (G1A-source).
\item[Gate 1B] Complementary signed/M\"obius/high-order Type-II provider based on
(G1B-center)--(G1B-M).
\item[Gate 2] Source-specific Ford--Maynard lower-bound compiler.
\item[BPP] Band-limited prime-participation mechanism converting moving-prime participation
to a family-energy saving in Gate~1A.
\item[SB-$\nu$] Synchronized two-view/source-glue fibre condition in the Gate-1A source
adapter.
\item[Pure5] Five-centred-defect distinguished-prime source family in the Gate-1B chain.
\item[NearPrim] Near-primitive/full-conductor slice ($r=1$) of the Gate-1B conductor family.
\item[QK56] Full same-/cross-modulus covariance parent (QK56-parent).
\item[SOURCE-ADAPTER45] Literal source dictionary from actual Ford-generated packets to
proved providers.
\item[$C_{\rm FM}$] Invocation-by-invocation certificate replacing the universal Type-II
hypothesis in the chosen Ford proof by generated-packet estimates.
\item[$R_9$] Balanced nine-coordinate Ford-generated leakage candidate; finite binomial
coefficient $70$ is proved, physical $G(d;n)$ specialization is not.
\item[FCL] Fixed-certificate leakage/correlation package (FCL).
\item[$N_2$] Exceptional Ford collar requiring a separate source-specific aggregate bound.
\item[Balanced7] Fixed-degree seven-prime-box V-cell at $Y=X^{1/7}$.
\item[V-cell] A generated packet in the model/factorial seven-prime branch.
\item[M\"obius singleton] Generated short Type-II factor carrying the M\"obius sign.
\item[fixed-$g_\ast$] One fixed positive Ford certificate used in the leakage transference.
\item[\texttt{WindowPairSupply}] The exact large-$M$ adjacent prime-power window predicate consumed by the finite
Lean compiler.
\end{description}


\section{The source-minimal affine algebra}

The current clean narrative starts from the exact identity $\Lambda=\mu*\log$.  The supplied
V15 Aristotle replay report records a successful $8118$-job kernel build of the declarations
summarized below.  Because public repository \texttt{main} has not yet synchronized that replay,
their status in this candidate is \status{LEAN-PROVED IN LATEST V15 REPLAY}, not
\status{PUBLIC LEAN-CHECKED}.  No analytic estimate is encoded by these identities.

\subsection{The identity $\Lambda=\mu*\log$}

For the genuine arithmetic functions,
\[
\mu*\log=\mu*(\Lambda*1)=(\mu*1)*\Lambda=\Lambda,
\]
and hence
\[
\boxed{\Lambda(N)=\sum_{qr=N}\mu(q)\log r.}
\]
At the affine value $N=2mn+s$, $s\in\{-1,+1\}$,
\[
\boxed{\Lambda(2mn+s)=\sum_{qr=2mn+s}\mu(q)\log r.}
\]
The V15 replay also preserves the natural-number sign firewall for the minus case.

\subsection{Exact factor partition and determinant-one line}

For a cutoff $U$, factor pairs split disjointly into
\[
q\le U,\qquad q>U,\ r\le U,\qquad q>U,\ r>U.
\]
The hard class is exact rather than heuristic:
\[
\sum_{\substack{qr=2mn+s\\q>U,\ r>U}}\mu(q)\log r.
\]
Every hard pair satisfies
\[
qr-2mn=s,\qquad s=\pm1.
\]
This implies
\[
(r,2m)=1,\qquad(q,2n)=1.
\]
Fixing $r,m$ and one solution $(q_0,n_0)$, all integer solutions are
\[
\boxed{q=q_0+2mt,\qquad n=n_0+rt,\qquad t\in\mathbb Z,}
\]
with unique $t$.  The V15 replay report records nonvanishing, the coprimality statements,
both directions of this parametrization, and multiplicity-one uniqueness as kernel-checked.

The same replay proves only the finite exponent ledger, not an asymptotic estimate.  In
particular,
\[
\frac16-\frac{16623}{100000}=\frac{131}{300000}>0.
\]

\section{The balanced-seven factorial Euler reduction}

Freeze the literal balanced-seven source
\[
Q=X^{3/5},\qquad Y=X^{1/7},
\]
and
\[
\B_{7,s}(X)=
\sum_{q\sim Q}\mu(q)
\left[
\sum_{\substack{p_1,\ldots,p_7\sim Y\\2p_1\cdots p_7\equiv-s\pmod q}}
\prod_{i=1}^7\omega_i(p_i)\,W(p_1\cdots p_7/X)
-\mathfrak M^{\mathrm{phys}}_{7,s}(q)
\right].
\]
For odd $q$, write
\[
 a_s(q)\equiv -s\,2^{-1}\pmod q.
\]
The target is $o(X/\log X)$.  The V15 Lean replay already proves the finite coefficient
\[
\sum_{j=0}^3(-1)^j\binom7j=-20
\]
and the earlier squarefree labelled polarization.  The controlling representation below is
newer and removes the collision remainder entirely.

\subsection{Factorial Euler polarization}

For $\mathbf z=(z_1,\ldots,z_7)\in\mathbb T^7$, define
\[
a_{\mathbf z}(p)=\frac17\sum_{i=1}^7z_i\omega_i(p)
\]
for primes in the actual box, and $a_{\mathbf z}(p)=0$ outside it.  Define a multiplicative
function $F_{\mathbf z}$ by the prime-power data
\[
\boxed{F_{\mathbf z}(p^e)=\frac{a_{\mathbf z}(p)^e}{e!},\qquad e\ge0.}
\]
This is ordinary multiplicativity; complete multiplicativity is neither required nor
desirable here.

Its local Dirichlet factor is
\[
\sum_{e\ge0}F_{\mathbf z}(p^e)p^{-es}
=\exp\!\left(a_{\mathbf z}(p)p^{-s}\right),
\]
so the generalized von Mangoldt coefficients satisfy
\[
\Lambda_{F_{\mathbf z}}(p)=a_{\mathbf z}(p)\log p,\qquad
\Lambda_{F_{\mathbf z}}(p^e)=0\quad(e\ge2).
\]
Since $|a_{\mathbf z}(p)|\le1$,
\[
|\Lambda_{F_{\mathbf z}}(n)|\le\Lambda(n),
\]
so $F_{\mathbf z}$ lies in the multiplicative class $\mathcal C$ used by
Drappeau--Granville--Shao.  This is a research/algebraic derivation; the factorial
polarization has not yet completed a Lean replay.

\subsection{Exact labelled coefficient extraction, including repeated primes}

Let
\[
n=\prod_pp^{e_p},\qquad\Omega(n)=\sum_pe_p=7.
\]
Then
\[
F_{\mathbf z}(n)=
\prod_p\frac1{e_p!}
\left(\frac17\sum_i z_i\omega_i(p)\right)^{e_p}.
\]
The factors $e_p!$ cancel exactly the multiplicity from permuting equal prime occurrences.
Therefore
\[
\boxed{
7^7[z_1\cdots z_7]F_{\mathbf z}(n)
=
\sum_{\substack{p_1\cdots p_7=n\\p_i\ {\rm ordered}}}
\prod_{i=1}^7\omega_i(p_i).
}
\]
Consequently,
\[
\boxed{
\begin{aligned}
&\sum_{\substack{p_1,\ldots,p_7\sim Y\\
2p_1\cdots p_7\equiv-s\pmod q}}
\prod_{i=1}^7\omega_i(p_i)\,W(p_1\cdots p_7/X)
\\
&\qquad =
7^7\int_{\mathbb T^7}\overline{z_1\cdots z_7}
\sum_{\substack{n\asymp X\\n\equiv a_s(q)\pmod q}}
F_{\mathbf z}(n)W(n/X)\,d\mathbf z .
\end{aligned}}
\]
This identity is exact for repeated primes.  Hence the former
$O(X^{6/7+o(1)})$ repeated-prime correction is retired from the controlling proof
architecture.  The earlier collision router remains a valid historical check in the
technical dossier.

\subsection{Expected-term linearity and the exact comparison pin}

Let $\mathcal E_q(f;a)$ be any expected-term operator that is linear in $f$:
principal projection, small-conductor subtraction, exceptional-character contribution, or
the literal combination adopted by the eventual analytic theorem.  Linearity gives
\[
\boxed{
7^7[z_1\cdots z_7]\mathcal E_q(F_{\mathbf z};a)
=
\mathcal E_q\!\left(7^7[z_1\cdots z_7]F_{\mathbf z};a\right).
}
\]
Thus polarization itself creates no comparison obstruction.

Define
\[
\mathfrak M^{\mathrm{fac}}_{7,s}(q)
=
7^7\int_{\mathbb T^7}\overline{z_1\cdots z_7}\,
\mathcal E_q(F_{\mathbf z};a_s(q))\,d\mathbf z
\]
and
\[
\Delta_{\mathcal E}(F_{\mathbf z};X,q,a)
=
\sum_{\substack{n\asymp X\\n\equiv a\pmod q}}
F_{\mathbf z}(n)W(n/X)-\mathcal E_q(F_{\mathbf z};a).
\]
Then
\[
\boxed{
\begin{aligned}
\B_{7,s}(X)
={}&7^7\int_{\mathbb T^7}\overline{z_1\cdots z_7}
\sum_{q\sim X^{3/5}}\mu(q)
\Delta_{\mathcal E}(F_{\mathbf z};X,q,a_s(q))\,d\mathbf z\\
&+\sum_{q\sim Q}\mu(q)
\bigl(\mathfrak M^{\mathrm{fac}}_{7,s}(q)
-\mathfrak M^{\mathrm{phys}}_{7,s}(q)\bigr).
\end{aligned}}
\]
There is no repeated-prime remainder.

The comparison problem is therefore not ``does polarization commute with the main term?''
It does.  The remaining source question is the physical identification
\[
\boxed{\mathfrak M^{\mathrm{fac}}_{7,s}(q)
=\mathfrak M^{\mathrm{phys}}_{7,s}(q)}
\]
at the required aggregate precision: principal local density together with the literal
exceptional-character convention.  This remains
\path{AFFINE287-MULOG-COMPARISON-LOWCOND-MATCH45}; it is not closed.

\subsection{The first exact V-branch analytic residual}\label{sec:factorial-endpoint}

The controlling analytic statement is now
\[
\boxed{
\begin{aligned}
\int_{\mathbb T^7}\overline{z_1\cdots z_7}
\sum_{q\sim X^{3/5}}\mu(q)\,
\Delta_{\mathcal E}
\left(F_{\mathbf z};X,q,a_s(q)\right)\,d\mathbf z
=o(X/\log X).
\end{aligned}}
\]
Programme label:
\[
\boxed{\begin{gathered}\text{\ttfamily AFFINE287-FACTORIAL-OMEGA7-}\\\text{\ttfamily SIGNED-ENDPOINT45}\end{gathered}}
\]
It is deliberately narrower than each of the following:
\begin{itemize}[leftmargin=1.6em]
\item a Bombieri--Vinogradov theorem for every $f\in\mathcal C$;
\item an absolute average $\sum_q|\Delta_q|$;
\item a theorem uniform in every $\mathbf z\in\mathbb T^7$;
\item a general smooth-number distribution theorem.
\end{itemize}
Only the signed $\mu(q)$-average of one degree-$(1,\ldots,1)$ torus coefficient is needed.

Together,
\[
\boxed{
\begin{array}{c}
\begin{gathered}\text{\ttfamily AFFINE287-FACTORIAL-OMEGA7-}\\\text{\ttfamily SIGNED-ENDPOINT45}\end{gathered}\\
+\ \texttt{AFFINE287-MULOG-COMPARISON-LOWCOND-MATCH45}\\[1mm]
\Downarrow\\
\texttt{AFFINE287-BALANCED7-MODULUS-AVERAGE45\ CLOSED}.
\end{array}}
\]
Both antecedents remain open.

\section{Published distribution results and the exact dictionary audit}

Drappeau--Granville--Shao work with the class $\mathcal C$ of
\emph{multiplicative} functions satisfying
$|\Lambda_f(n)|\le\Lambda(n)$; complete multiplicativity is only a special case.
Their theorem reaches moduli $x^{3/5-\varepsilon}$ for sufficiently smooth support, so it
does not literally supply the exact endpoint $Q=X^{3/5}$ required above.

For Pascadi's current arXiv v2, the relevant stated result is the smooth-number theorem
(Theorem 1.5, refined as Theorem 6.4).  Immediately after Theorem 1.5 the paper remarks
that a similar result for smooth-supported multiplicative functions can be deduced by a
slight extension of Proposition 6.3; the current v2 does not state a separate
``Theorem 1.7'' giving the black box required here.  More importantly, the written proof
architecture itself is quantitatively nonclosing at the seven-prime smoothness scale.

Indeed, renaming the proof's internal small parameter as $\eta$, the proof arranges
\[
Q\le x^{5/8-100\eta},\qquad y\le x^\eta.
\]
At the required modulus
\[
Q=x^{3/5},
\]
the first inequality implies
\[
\frac35\le\frac58-100\eta
\quad\Longrightarrow\quad
\boxed{\eta\le\frac1{4000}}.
\]
But the source has
\[
y=x^{1/7},
\]
so the second inequality would require
\[
\boxed{\eta\ge\frac17}.
\]
These conditions are incompatible.  Thus the \emph{direct existing Pascadi proof route},
with its written parameters, does not furnish the factorial endpoint theorem at
$(Q,y)=(X^{3/5},X^{1/7})$.  This is a theorem-dictionary/proof-parameter audit, not a
claim that no future extension of Pascadi's technology can succeed.

This explains the current research direction: preserve the literal fixed degree seven rather than forcing a
generic smooth-factorization parameter, carry that structure through dispersion, and test a post-dispersion
four-block arrangement.  The ``fixed-degree-seven / four-block'' programme is a research direction, not a
proved theorem or an already matched Gate-1B provider.  Its purpose is precisely to attack the narrower signed
torus coefficient without paying for a theorem uniform over all smooth-supported functions.

The earlier V14 Vaughan decomposition remains mathematically valid as a nonminimal
alternative source audit.  It is no longer presented as a co-equal controlling branch in
this public narrative; its prime-modulus two-outer packet and related failed dictionaries
belong in the technical provenance appendix.

\section{The current end-to-end dependency graph}

The unconditional finite theorem and the conditional large-$M$ interface remain the two endpoints of the
project:
\[
\boxed{\begin{array}{c}
\text{exact public problem}\\
\downarrow\\
\text{prime-power window compiler + finite certificates}\\
\downarrow\\
M\le4\times10^9\text{ excluded}
\end{array}}
\]
while the conditional completion is
\[
\boxed{\begin{array}{c}
\text{eventual }\WPS(M)\text{ supply}\\
\downarrow\\
\text{public Lean closure compiler}\\
\downarrow\\
\text{Erd\H{o}s \#287.}
\end{array}}.
\]

Between them sits the restored analytic gate architecture:
\[
\boxed{
\begin{array}{c}
\text{Gate 0: physical sequence class, comparison and Type I}\\
\downarrow\\
\text{Ford-generated packet census / source adapter}\\
\downarrow\\
\text{Gate 1A}\quad\vert\quad\text{Gate 1B}\quad\vert\quad\text{direct providers}\\
\downarrow\\
\text{generated-(7.23) or universal Full-Type-II reassembly }+C_{\rm FM}\\
\downarrow\\
\text{Gate 2: physical }H(n),\ N_2,\ \varepsilon_*\text{ and lower-bound geometry}\\
\downarrow\\
\text{positive affine mass and effectivity}\\
\downarrow\\
\WPS(M)\text{ or a stronger sign-sensitive witness.}
\end{array}}
\]

The current 287 overlays enter this spine at the generated-packet level:
\[
\Lambda=\mu*\log,\qquad
F_{\mathbf z}(p^e)=a_{\mathbf z}(p)^e/e!,\qquad
\texttt{FACTORIAL-OMEGA7-SIGNED-ENDPOINT45},
\]
together with the physical comparison identification.  Even if the balanced-seven pair of residuals is closed,
one still needs the remaining V-cells/M-branch, the fixed-$g_*$ packet census, provider assignment, balanced
$R_9$ source specialization where relevant, FCL/leakage reassembly, $N_2$, the Gate-2 lower bound and an
effective large-$M$ witness.  This is why the clean paper retains the Gate spine rather than compressing every
analytic arrow into the phrase ``hard cancellation''.

\section{Formal verification and reproducibility}

\subsection{Public Lean bank}

Public repository \texttt{main} presently contains the exact problem predicate, the
finite exclusion through $4\times10^9$, the \texttt{WindowPairSupply} closure compiler,
the V14 structural source bank, and at least some individual V15 source modules such as
\path{AffineMuLogHardSource.lean}.  The aggregate \path{Main.lean} import spine and V15
status file are not yet synchronized.  The claims labelled \status{PUBLIC LEAN-CHECKED}
are therefore limited to declarations directly verified on public \texttt{main}; the full V15 replay status
below is tracked separately.

\subsection{Latest supplied V15 Aristotle replay}

A newer user-supplied Aristotle archive and build report contain the following V15 files:
\begin{itemize}[leftmargin=1.7em,itemsep=0.15em]
\item \path{AffineMuLogIdentity.lean};
\item \path{AffineMuLogHardSource.lean};
\item \path{AffineMuLogLine.lean};
\item \path{AffineMuLogExponentLedger.lean};
\item \path{BalancedSevenFinite.lean};
\item \path{BalancedSevenPolarization.lean};
\item \path{Erdos287V15Status.lean}.
\end{itemize}
The report states:
\begin{itemize}[leftmargin=1.6em]
\item \texttt{lake build} succeeded with $8118$ jobs and no errors;
\item a Lean-code scan found no \texttt{sorry}, \texttt{admit}, user
  \texttt{axiom}, \texttt{opaque}, \texttt{unsafe}, \texttt{native\_decide}, or
  \texttt{@[implemented\_by]};
\item \texttt{\#print axioms} for the principal new theorems used only subsets of
  \texttt{propext}, \texttt{Classical.choice}, and \texttt{Quot.sound};
\item the analytic structures \path{PolarizedOmega7SignedEoD} and
  \path{MuLogComparisonLowCondMatch} were deliberately left uninhabited.
\end{itemize}
Accordingly this candidate labels the exact V15 algebra
\status{LEAN-PROVED IN LATEST V15 REPLAY}.  At the time of this editorial audit,
public \texttt{main} still did not display that full replay, so public synchronization
is kept as a separate provenance fact rather than silently conflated with kernel status.

\subsection{Factorial Euler layer}

The latest supplied V16 factorial Aristotle replay reports the factorial Euler polarization,
exact repeated-prime encoding, formal local Euler-series/generalized-$\Lambda$ identities,
expected-term linearity, and the Pascadi parameter ledger as kernel-checked.  The build report
states $8123$ jobs, zero errors, no unfinished-proof or unsafe-code markers, and principal axiom
reports restricted to the standard logical principles \texttt{propext},
\texttt{Classical.choice} and \texttt{Quot.sound}.  Because the exact factorial replay has not
been independently synchronized into every public aggregate status file, this paper labels it
\status{LEAN-PROVED IN LATEST V16 FACTORIAL REPLAY} rather than silently equating replay status
with public-main synchronization.

\subsection{Prepublication state}

This R4 candidate preserves the clean front half, inserts a self-contained Gate 0--2 theory and verification manual, and retains the organized current frontier ledger.  It is not a GitHub publication candidate until the current Gate-1B research ledger and any newer Aristotle replay are re-audited.

This file is a local prepublication candidate.  It has not been pushed to the canonical
GitHub PDF or the legacy forum PDF path.  The existing public V16 and the immutable V15
archive remain untouched while this candidate is reviewed.

\section{Conclusion}

The strongest unconditional result in the project remains finite: there is no exact
counterexample to Erd\H{o}s Problem \#287 with maximum denominator at most
$4\times10^9$.  The exact global remaining interface is equally clear: establish effective
eventual \texttt{WindowPairSupply}.

On the analytic side, the source-minimal $\Lambda=\mu*\log$ algebra and determinant-one
line have now been reported kernel-checked in the latest V15 replay.  For the balanced-seven
cell, the controlling representation is no longer the squarefree polarization plus a
collision remainder.  The factorial Euler function
\[
F_{\mathbf z}(p^e)=a_{\mathbf z}(p)^e/e!
\]
recovers the ordered seven labelled source exactly, including repeated primes, and
polarization commutes with every linear expected-term operator.

The first exact V-branch analytic open is therefore
\[
\boxed{\begin{gathered}\text{\ttfamily AFFINE287-FACTORIAL-OMEGA7-}\\\text{\ttfamily SIGNED-ENDPOINT45}\end{gathered}},
\]
together with the source identification
\[
\boxed{\texttt{AFFINE287-MULOG-COMPARISON-LOWCOND-MATCH45}}.
\]
The endpoint theorem is much narrower than a general class-$\mathcal C$ Bombieri--Vinogradov
statement.  The current DGS and Pascadi theorems do not literally furnish it, and the
direct written Pascadi proof parameters are incompatible with
$(Q,y)=(X^{3/5},X^{1/7})$.

Closing those two balanced-seven inputs would still leave the remaining V-cells and M-branch, the literal
Ford-generated packet census and provider assignment, the $k=0$ smooth packet, any balanced-$R_9$ physical
$G(d;n)$ specialization actually generated by the sieve, fixed-certificate leakage/FCL, the source-specific
$N_2$ and Gate-2 lower-bound compiler, effectivity, and eventual large-$M$ supply.  Gate 1A is required only if
the literal census routes a needed packet to its common-weight provider.  Erd\H{o}s Problem \#287 therefore
remains open.


\clearpage
\part{Current Frontier and Provenance Ledger}
\label{part:technical-ledger}

\section{Purpose, coverage baseline and status firewall}
\label{sec:technical-baseline}

Part~\ref{part:technical-ledger} restores the \emph{technical depth} of the earlier public
dossier without restoring its chronological patch-stack format.  Its coverage floor is the
V15 public-review snapshot committed as
\[
\boxed{\texttt{7fc3797105519c4e6a1669f174ce7394336bdc18}},
\]
which published the former 50-page-style Gate-1A/Gate-1B/Ford programme together with the
then-current balanced-seven update.  That snapshot is used here only as a coverage baseline:
when later work changes the controlling route, the later status wins.

Three time/status layers are kept separate throughout:
\begin{enumerate}[leftmargin=1.7em]
\item \textbf{public/formal bank}: theorem or compiler actually present in the public Lean
  project;
\item \textbf{latest supplied research ledger}: a mathematically audited frontier supplied
  after the V15 snapshot but not silently promoted to a public Lean or published theorem;
\item \textbf{open/source interfaces}: deliberately unproved analytic or source-identification
  obligations.
\end{enumerate}

This is important for Gate~1B.  The public formal socket still records the one-completion
coefficient
\[
\beta_{D,P}(d,p)=\mu_D(d)\Lambda_P(p)
\]
and the exact $1/104$ distribution-margin budget.  That socket is neither false nor erased.
However, the latest supplied analytic work has moved beyond it to a Motohashi/rank-one
frontier.  The formal socket is therefore a \emph{banked abstraction/alternative provider
interface}, while the rank-one chain below is the current research frontier.

Detailed dead-route archaeology remains outside the main narrative.  What is restored here
is every current component needed to reconstruct the proof programme: Gate~1A source class,
Gate~1B current frontier and downstream firewall, the Ford-generated packet census,
balanced-$R_9$/fixed-certificate leakage, Gate~2, and the relation of the current
factorial-$\Omega_7$ source to those providers.

\section{Gate 1A: current canonical-provider ledger}
\label{sec:gate1a-ledger}

Gate~1A is a conditional theorem for one canonical common-weight physical operator.  It is
not a universal theorem for arbitrary row-dependent source weights and it is not assumed
globally necessary for Erd\H{o}s~\#287.

For a row
\[
e=(r,k,m),\qquad m'=m+kr,
\]
the affine forms satisfy
\[
m'A_e(t)-mB_e(t)=2k,
\]
and the canonical physical operator is
\[
C_e(b,d)=
\sum_tW_D(t)
\left(\sum_{p\asymp L}b_p\rho_p(A_e(t))\right)
\left(\sum_{q\asymp L}d_q\rho_q(B_e(t))\right).
\]
The conditional target remains
\[
\sum_{e\in\mathcal E^\sharp}|C_e(b,d)|^2
\ll
\frac{ML^4}{H}X_1^{o(1)}
\sum_{p\asymp L}|b_p|^2
\sum_{q\asymp L}|d_q|^2,
\]
with the same normalization conventions as in Part~I.

\subsection{What is banked and what is still external}

\begin{center}
\renewcommand{\arraystretch}{1.18}
\begin{tabularx}{\textwidth}{>{\raggedright\arraybackslash}p{0.31\textwidth}
>{\raggedright\arraybackslash}p{0.20\textwidth}X}
\toprule
Gate-1A component & Status & Current interpretation\\
\midrule
Determinant/row algebra and all-$m$ row family & \status{BANKED FINITE/OPERATOR} &
The theorem is not restricted to clean-$P_3$ rows.\\
Canonical common physical weight $W_D$ & \status{SOURCE FIREWALL BANKED} &
Later row dependence must be deterministically derived from this one datum.\\
Scale ledger, row-energy conservation and normalization bridge & \status{BANKED} &
Includes the $H^2$ conversion and the signed vertex margins.\\
Finite BPP participation-to-energy compiler & \status{BANKED} &
Participation implies a moving-family energy saving once the physical participation input is supplied.\\
Projective ratio-class/axis algebra & \status{BANKED FINITE ALGEBRA} &
The state diagonal is routed inside the zero-curvature/projective branch rather than discarded.\\
Outer-square-root algebra and exponent budget & \status{BANKED} &
The formal compiler checks the square-root conversion but does not manufacture the physical energy pin.\\
Poisson--Bruhat transport/source ceiling & \status{SOURCE/ANALYTIC PIN} &
Must preserve the literal physical packet and absolute comparator.\\
Fixed-state exclusion and prime participation & \status{EXTERNAL/SOURCE PIN} &
The finite compiler is present; the physical smooth-envelope/short-interval input is external.\\
Outer-root physical energy identification & \status{SOURCE OPEN} &
Identifies the abstract BPP energy with the actual Gate energy.\\
Projective/diagonal physical source routing & \status{SOURCE OPEN} &
Finite algebra exists; the physical source ceiling must be supplied at target scale.\\
Literal recombination estimate & \status{SOURCE OPEN} &
Controls discarded/excision/recombination errors without spending the same saving twice.\\
Arbitrary edge-dependent $D_2$ family & \status{OUTSIDE CANONICAL THEOREM} &
Requires an upstream Gate-0 adapter or finite-template/nuclear reduction.\\
\bottomrule
\end{tabularx}
\end{center}

The current source-glue issue can be summarized as a synchronized two-view/fibre problem:
the same physical packet has to satisfy the common-weight representation, the BPP
participation representation and the final source-routing representation simultaneously.
We refer to this compactly as the \emph{SB-$\nu$/source-glue interface}.  It is a
source-identification problem, not an exponent deficit in the already-banked finite
operator algebra.

\subsection{When Gate 1A is required for \#287}

The correct logical statement is
\[
\boxed{
\texttt{GATE1A\_REQUIRED}
\ \text{is decided by the literal Ford-generated packet census.}
}
\]
An actual generated packet may use Gate~1A only after a scale/source dictionary supplies:
\begin{itemize}[leftmargin=1.6em]
\item the relation between the Gate ambient scale $X_1$ and the \#287 scale $X$;
\item the identification of the physical bilinear variables with the Gate row variables;
\item the induced $R,L,H,K,U,D$ ranges;
\item preservation of the one common physical weight, or a proved finite-template/nuclear
  reduction with acceptable total cost;
\item survival of the binding fixed-power margin after truncation, multiplicity and
  scale conversion.
\end{itemize}
No global Gate-1A closure is assumed if the generated census never routes a needed packet
to this source class.

\section{Gate 1B: from the banked one-completion socket to the current rank-one frontier}
\label{sec:gate1b-ledger}

Gate~1B is the complementary signed/parity-breaking provider.  Its defining feature is
that the M\"obius sign and shifted congruence data must remain attached to the physical
source; squaring or replacing the source by an arbitrary coefficient family can destroy
the resource that is meant to break parity.

\subsection{Banked one-completion socket}

The public challenge file records the authoritative one-completion coefficient
\[
\boxed{\beta_{D,P}(d,p)=\mu_D(d)\Lambda_P(p)}
\]
as separately linear in the M\"obius and prime variables.  It also proves the exact dyadic
budget
\[
\theta_{\rm req}=\frac12+\frac1{104}
\]
and the implication
\[
\begin{gathered}
\text{literal source-to-convolution-BV interface}\\
\text{at level }\theta\ge\theta_{\rm req}\\
\Downarrow\\
\text{abstract Gate-1B convolution conclusion}
\end{gathered}
\]
The source-to-interface theorem is deliberately open.  This formal socket remains useful,
but it no longer describes the first research-level analytic obstruction.

\subsection{Latest supplied Motohashi/rank-one chain}

The latest Gate-1B research ledger supplied for this candidate gives the following
controlling progression.  These labels are \emph{research-status labels}, not claims of
public Lean or published theorems.

\begin{center}
\renewcommand{\arraystretch}{1.17}
\begin{tabularx}{\textwidth}{>{\ttfamily\raggedright\arraybackslash}p{0.42\textwidth}
>{\raggedright\arraybackslash}p{0.22\textwidth}X}
\toprule
Node & Latest status & Role\\
\midrule
MOTOHASHI-ABC-EXACT-PIN45 & \status{PASS} &
Exact A/B/C source pin for the Motohashi decomposition.\\
TWISTED-DEFECT-ABC45 & \status{PASS} &
Carries the centred/twisted defect data through the A/B/C split.\\
FIVEFOLD-MOTOHASHI-ITERATION45 & \status{PASS} &
Iterates the exact five-defect source without replacing it by unrelated coefficient norms.\\
RANKONE-POLYLOGK-INTERIOR45 & \status{CLOSED MODULO COMPARISON} &
Fixed-polylogarithmic-$k$ interior reduced to the comparison/main-term convention.\\
RANKONE-ENDPOINT-U-DIAGONAL45 & \status{POWER CLOSED} &
Endpoint $u$-diagonal carries a saving $X^{-1/12+o(1)}$.\\
RANKONE-ENDPOINT-U-OFFDIAG45 & \status{FIRST ANALYTIC OPEN} &
Current first analytic frontier in the rank-one endpoint family.\\
RANKONE-HIGHK45 & \status{OPEN} &
High-$k$ regime is not covered by the fixed-polylogarithmic-$k$ interior analysis.\\
PURE5-COMPARISON-MAINTERM-PIN & \status{SOURCE / COMP. OPEN} &
Literal physical main-term normalization must match the pure-five source.\\
\bottomrule
\end{tabularx}
\end{center}

Thus the present first research question is no longer the abstract $1/104$ socket and no
longer a broad ``full covariance'' theorem.  It is the source-preserving endpoint
off-diagonal problem
\[
\boxed{\begin{gathered}\text{\ttfamily RANKONE-ENDPOINT-}\\\text{\ttfamily U-OFFDIAG45}\end{gathered}}
\]
The comparison pin is logically parallel: closing the off-diagonal estimate without the
physical main-term convention does not close the pure-five packet.

\subsection{Downstream Gate-1B firewall}

Even a solution of the first endpoint off-diagonal does not jump directly to Gate~1B.
The current source ledger keeps the following downstream structure explicit:
\[
\boxed{
\begin{array}{c}
\text{Motohashi A/B/C + twisted fivefold source}\\
\downarrow\\
\text{fixed-polylog }k\text{ interior}
\quad+\quad
u\text{-diagonal endpoint}\\
\downarrow\\
\boxed{\begin{gathered}\text{\ttfamily RANKONE-ENDPOINT-}\\\text{\ttfamily U-OFFDIAG45}\end{gathered}}\\
\parallel\\
\text{\ttfamily PURE5-COMPARISON-MAINTERM-PIN}\\
\downarrow\\
\text{high-}k\text{ sector}\\
\downarrow\\
\text{\ttfamily PURE5-DP-SIGNED45}\\
\downarrow\\
|J|=4,3,2,1\text{ defect strata}\\
\downarrow\\
\text{\ttfamily NEARPRIM}\\
\downarrow\\
r>1\text{ complement}+\text{transition}+\text{recursion}\\
\downarrow\\
\text{\ttfamily QK56-FULL-PARENT}\\
\downarrow\\
\text{\ttfamily QK56-EXHAUSTIVENESS}\\
\downarrow\\
\text{shifted }TT^\ast\text{ source/reassembly}\\
\downarrow\\
\texttt{GATE1B}.
\end{array}}
\]
This diagram is a \emph{dependency firewall}: later nodes are not declared solved merely
because an earlier parent is controlled.

\subsection{Where QK56 sits}

The QK56 parent is a full covariance object of the form
\[
\mathcal C_{\rm QK}
=
\sum_{q_1,q_2}
\sum_{h_1,k_1,h_2,k_2}
A(q_1,q_2;h_1,k_1,h_2,k_2)
Q(q_1;h_1,k_1)\,Q(q_2;h_2,k_2).
\]
The split $q_1=q_2$ versus $q_1\ne q_2$ is the same-modulus/cross-modulus split of one
parent, not two unrelated theorems.  Earlier work established substantial same-$q$ and
cross-$q$ children and the canonical zero mode at the finite/algebraic level.  The present
rank-one chain places QK56 later in the exhaustiveness/reassembly firewall rather than at
the first analytic open.

The exact status therefore is:
\[
\boxed{
\begin{gathered}
\text{QK56 remains structurally live, but}\\
\begin{gathered}\text{\ttfamily RANKONE-ENDPOINT-}\\\text{\ttfamily U-OFFDIAG45}\end{gathered}\\
\text{is earlier in the current closure order}
\end{gathered}}
\]

\section{Ford--Maynard generated packets: census before provider}
\label{sec:ford-census}

The Ford--Maynard universal Type-II hypothesis is a clean published sufficient wrapper,
but Proposition~7.22 does not invoke arbitrary coefficient pairs out of thin air.  After a
finite decomposition it generates the grouped equation-(7.23) family
\[
\sum_{n=n'n''\asymp x}Y_1(n')Y_2(n'')w_{n'n''}.
\]
A proof-specific replacement is therefore legitimate only after two source obligations:
\begin{enumerate}[leftmargin=1.7em]
\item a \emph{literal generated-packet census} listing every coefficient family produced by
  the proof;
\item a replacement certificate $C_{\rm FM}$ proving that every Type-II invocation is
  covered by one proved provider after all Mellin/Perron, support and norm transformations.
\end{enumerate}

The programme labels
\[
\begin{gathered}
\text{\ttfamily SOURCE-ADAPTER45},\\
\text{\ttfamily FM-TO-GATE-COORDINATE-CENSUS45}
\end{gathered}
\]
refer to this source dictionary.  Neither can be replaced by the slogan ``the packet looks
Type-II''.

\subsection{Current provider-facing census}

\begin{center}
\renewcommand{\arraystretch}{1.18}
\begin{tabularx}{\textwidth}{>{\raggedright\arraybackslash}p{0.22\textwidth}
>{\raggedright\arraybackslash}p{0.23\textwidth}
>{\raggedright\arraybackslash}p{0.22\textwidth}X}
\toprule
Generated family & Literal status & Candidate provider & Current blocker\\
\midrule
M\"obius singleton & source-specific generated packet &
Gate~1B / signed source &
literal dictionary and stability under generated outer factors\\
Model singleton & source-specific generated packet &
QK56/direct affine completion &
native packet dictionary; short M\"obius carrier may sit in the complement\\
Balanced-seven V-cell & explicit fixed-degree source &
factorial-$\Omega_7$ endpoint theorem &
signed endpoint + physical comparison identification\\
Balanced $R_9$ leakage & Ford-generated leakage candidate &
fixed-certificate/FCL lane &
literal $G(d;n)$ specialization and source reassembly\\
Common-weight/edge packets & census-dependent &
Gate~1A only after exact match &
scale/source adapter and common-weight preservation\\
Pure-five/high-order packets & Gate-1B source family &
Motohashi/rank-one chain &
endpoint off-diagonal, comparison, high-$k$, downstream strata\\
Remaining fixed-$g_\ast$ packets & not exhaustively censused &
TBD from literal packet shape &
finite census and provider assignment\\
\bottomrule
\end{tabularx}
\end{center}

This table encodes the main logical correction:
\[
\boxed{
\text{provider assignment follows the generated source; it is not imposed globally.}
}
\]
In particular, Gate~1A is neither deleted nor assumed universally necessary.

\section{Balanced $R_9$, fixed-certificate leakage and the Gate-2 landing}
\label{sec:r9-fcl}

The balanced order-nine identity
\[
\sum_{j=0}^{4}(-1)^j\binom9j=70
\]
is a genuine finite/binomial fact.  The formal bank also records the complementary
alternating identity.  What is \emph{not} automatically proved is that Ford's actual
arithmetic weight $G(d;n)$ selects exactly this divisor-order specialization on the physical
balanced-$R_9$ cell.

Therefore the correct status is:
\[
\boxed{
\text{balanced }R_9
=
\text{Ford-generated leakage candidate,}
}
\]
with
\[
\boxed{
H(n)=70
\quad\text{only after the literal }G(d;n)\text{ specialization is proved.}
}
\]
The old $R_9$ death certificate and the abandoned unbalanced-$1/5$ redesign are not
revived.

The fixed-certificate architecture uses one explicit positive Ford certificate $g_\ast$
and its arithmetic weight
\[
H_\ast(n)=\sum_{d\mid n}G_\ast(d;n).
\]
Its transference identity separates:
\begin{itemize}[leftmargin=1.6em]
\item total correlation;
\item leakage outside the prime/main-composite regions;
\item the exceptional $N_2$ collar;
\item the positive comparison margin.
\end{itemize}
Balanced $R_9$ is one generated leakage family that can feed this architecture once the
physical source dictionary is proved; it is not the whole fixed-certificate leakage
problem.

\subsection{Gate 2 consumes, rather than invents, the packet bounds}

Gate~2 is the Ford--Maynard lower-bound/endgame compiler.  It fixes one compatible
$\varepsilon_\ast$ and the shrunk parameter triple
\[
P_{\varepsilon_\ast}
=
\left(
\frac12-\varepsilon_\ast,\,
\varepsilon_\ast,\,
\frac16-2\varepsilon_\ast
\right).
\]
Its source-specific replay requires:
\begin{enumerate}[leftmargin=1.7em]
\item Gate-0 comparison regularity, growth and exact Type-I input;
\item the literal $G(d;n)/H(n)$ dictionary and the Ford-generated packet census;
\item the source-specific aggregate $N_2$ estimate;
\item one compatible fixed $\varepsilon_\ast$;
\item either the universal Type-II wrapper, or all generated-(7.23) packet estimates
  together with \path{SOURCE-ADAPTER45} and $C_{\rm FM}$.
\end{enumerate}
Thus the correct closure order is
\[
\text{packet census}
\to
\text{Gate providers}
\to
\text{generated Type-II reassembly}
\to
\text{Gate 2 lower bound},
\]
not Gate~2 followed by a later packet census.

\section{Current \#287 fixed-degree-seven direction inside the gate programme}
\label{sec:fixed-degree-seven-ledger}

Part~I gives the controlling factorial Euler encoding
\[
F_{\mathbf z}(p^e)=\frac{a_{\mathbf z}(p)^e}{e!},
\qquad
a_{\mathbf z}(p)=\frac17\sum_{i=1}^{7}z_i\omega_i(p),
\]
which exactly includes repeated primes.  The first V-branch analytic residual is
\[
\boxed{\begin{gathered}\text{\ttfamily AFFINE287-FACTORIAL-OMEGA7-}\\\text{\ttfamily SIGNED-ENDPOINT45}\end{gathered}}
\]
and the physical comparison identification
\[
\boxed{\texttt{AFFINE287-MULOG-COMPARISON-LOWCOND-MATCH45}}
\]
remains separate.

The direct equal-seven three-block dictionary into Pascadi Proposition~6.3 fails its range
geometry, and the direct smooth-number/multiplicative-function proof architecture is
nonclosing at $y=X^{1/7}$ for the exact parameter reason proved in Part~I.  The current
research direction is therefore \emph{post-dispersion fixed-degree seven}: keep the degree
seven source literal, perform the relevant dispersion/completion first, and search for a
four-block or comparable structured theorem only after the physical source has exposed the
correct variables.  This is a research direction, not a theorem or an already-identified
published provider.

If the factorial endpoint and comparison pin are both proved, the balanced-seven cell
closes.  That implication is deliberately local:
\[
\boxed{
\text{Balanced7 closed}
\not\Longrightarrow
\text{FCL closed}
\not\Longrightarrow
\texttt{WindowPairSupply}.
}
\]
The remaining packet census and Gate-provider assignment still intervene.

\section{Master current-status ledger}
\label{sec:master-ledger}

{\footnotesize\begin{longtable}{>{\raggedright\arraybackslash}p{0.32\textwidth}>{\raggedright\arraybackslash}p{0.22\textwidth}>{\raggedright\arraybackslash}p{0.33\textwidth}}
\toprule
Node & Current status & Publication interpretation\\
\midrule
\endfirsthead
\toprule
Node & Current status & Publication interpretation\\
\midrule
\endhead
Finite exclusion $M\le4\times10^9$ & \status{PUBLIC LEAN-CHECKED} &
strongest unconditional theorem\\
\texttt{WindowPairSupply} compiler & \status{PUBLIC LEAN-CHECKED / CONDITIONAL} &
exact global large-$M$ interface\\
Gate 0 comparison/growth framework & \status{MIXED PASS / SOURCE PINS} &
literal Type-I/publication dictionary still matters\\
Gate 1A finite/operator mechanism & \status{SUBSTANTIALLY BANKED} &
conditional canonical common-weight provider\\
Gate 1A physical five-pin/source glue & \status{OPEN SOURCE/ANALYTIC} &
needed only for packets censused into Gate~1A\\
Gate 1B one-completion $1/104$ socket & \status{BANKED FORMAL INTERFACE} &
valid abstraction, not first research frontier\\
Motohashi A/B/C and twisted fivefold source & \status{LATEST RESEARCH PASS} &
current Gate-1B source reduction\\
Rank-one polylog-$k$ interior & \status{CLOSED MODULO COMPARISON} &
latest supplied research ledger\\
Rank-one $u$-diagonal endpoint & \status{POWER CLOSED }$X^{-1/12+o(1)}$ &
latest supplied research ledger\\
\path{RANKONE-ENDPOINT-U-OFFDIAG45} & \status{FIRST GATE-1B ANALYTIC OPEN} &
current rank-one frontier\\
\path{RANKONE-HIGHK45} & \status{OPEN} &
separate high-$k$ regime\\
\path{PURE5-COMPARISON-MAINTERM-PIN} & \status{SOURCE OPEN} &
parallel normalization requirement\\
Pure-five $\to$ defect strata $\to$ near-primitive chain & \status{DOWNSTREAM OPEN/ROUTING} &
must not be skipped after rank-one closure\\
QK56 full parent/exhaustiveness & \status{LATER GATE-1B FIREWALL} &
structurally live, not first current open\\
Ford generated-(7.23) census & \status{OPEN SOURCE/FINITE CENSUS} &
decides provider assignment\\
\path{SOURCE-ADAPTER45}, $C_{\rm FM}$ & \status{OPEN} &
required for proof-specific replacement of universal Type II\\
Balanced $R_9$ binomial $70$ & \status{FINITE PASS} &
physical $G(d;n)$ specialization still source open\\
Factorial Euler $\Omega_7$ encoding & \status{RESEARCH PASS / LEAN PENDING} &
current balanced-seven representation\\
Factorial signed endpoint & \status{OPEN ANALYTIC} &
current exact V-branch residual\\
Factorial/physical comparison match & \status{OPEN SOURCE} &
principal local density + literal exceptional-character convention\\
Remaining fixed-$g_\ast$ census & \status{OPEN} &
balanced-seven is only one generated family\\
FCL / positive affine mass & \status{OPEN} &
requires all leakage/provider inputs\\
Gate 2 lower-bound replay & \status{CONDITIONAL / NOT ACTIVATED} &
consumes generated packet bounds, $N_2$, $G/H$, $\varepsilon_\ast$\\
Eventual large-$M$ supply/effectivity & \status{OPEN} &
needed before the finite Lean compiler applies globally\\
Erd\H{o}s Problem \#287 & \status{OPEN} &
no unconditional proof claim\\
\bottomrule
\end{longtable}}

\section{Superseded, false and provenance-only routes}
\label{sec:provenance-only}

The following are preserved for auditability but are not allowed to displace the current
frontier:
\begin{itemize}[leftmargin=1.6em]
\item the V15 squarefree polarization plus $X^{6/7+o(1)}$ repeated-prime router:
  valid historical representation, superseded as controlling by the factorial Euler encoding;
\item complete multiplicativity as a required repair: \emph{false}; ordinary multiplicativity
  in class $\mathcal C$ is enough for the relevant dictionary;
\item a stated Pascadi ``Theorem 1.7'' multiplicative-function black box:
  not present in the audited current arXiv v2;
\item the direct Pascadi v2 proof route at $(Q,y)=(X^{3/5},X^{1/7})$:
  audited nonclosing because $\eta\le1/4000$ conflicts with $\eta\ge1/7$;
\item the V14 Vaughan/prime-modulus M\"obius two-outer route:
  mathematically valid but nonminimal, retained as source audit/provenance;
\item direct singleton $=$ existing Gate-1B common-conductor covariance:
  false as a literal source dictionary;
\item direct canonical Gate-1A determinant adapter for the affine determinant-one source:
  fails literally ($2k$ versus $\pm1$);
\item equal-seven Pascadi Proposition~6.3 three-block dictionary:
  false at the literal equal-seven geometry;
\item ordinary separated multiplicative large sieve / separated high moments:
  nonclosing by the previously audited $X^{1/10}$ deficit;
\item modulus micro-switch as an exhaustive closure:
  nonexhaustive because the prime-modulus sector survives.
\end{itemize}

The public paper therefore retains V15's \emph{diagnostic density} while keeping obsolete
routes visually and logically subordinate to the current closure chain.


\appendix
\section{Formal theorem and provenance map}

\subsection*{Public Lean bank}
\begin{itemize}[leftmargin=1.6em,itemsep=0.4em]
\item \path{RequestProject/Erdos287/ProblemStatement.lean}: exact problem predicate and
  ordered public-statement compiler.
\item Finite-range files: adjacent-hole obstruction, prime-power window blocker, and
  \path{no_Erdos287Counterexample_of_max_le_4e9}.
\item \path{RequestProject/Erdos287/ClosureInputs.lean}: \path{WindowPairSupply} and
  the end-to-end conditional closure compiler.
\item V14 source files/status: valid Vaughan structural audit, retained as provenance rather
  than the controlling source-minimal route.
\end{itemize}

\subsection*{Lean-proved in the latest supplied V15 replay}
The supplied V15 build report records kernel-checked proofs of:
\begin{itemize}[leftmargin=1.6em,itemsep=0.35em]
\item $\Lambda=\mu*\log$ and its affine specialization;
\item exhaustive small-$q$/small-$r$/hard recombination;
\item the determinant-one line, coprimality, iff parametrization, and uniqueness;
\item the finite exponent ledger;
\item the balanced-seven binomial coefficient $-20$;
\item the repeated-prime finite router;
\item the squarefree labelled polarization and its $7^{-7}$ normalization.
\end{itemize}
The same status source leaves the analytic endpoint and comparison structures uninhabited.

\subsection*{Research-proved / Lean pending}
The current hostile audit adds:
\begin{itemize}[leftmargin=1.6em,itemsep=0.35em]
\item factorial Euler polarization
  $F_{\mathbf z}(p^e)=a_{\mathbf z}(p)^e/e!$;
\item exact ordered coefficient extraction including repeated primes;
\item $F_{\mathbf z}\in\mathcal C$ from its Euler logarithmic derivative;
\item exact commutation with every linear expected-term operator;
\item the proof-parameter obstruction
  $\eta\le1/4000$ versus $\eta\ge1/7$ for the direct Pascadi v2 route.
\end{itemize}
The latest supplied V16 factorial replay is claimed as kernel-checked; public-main synchronization remains separate.

\subsection*{Gate and Ford--Maynard source provenance}
The Gate architecture in Section~\ref{sec:theory-purpose} is a compressed restoration of the earlier
standalone verification dossier, not a reconstruction from theorem names alone.  The latest supplied formal
bank further records: the Gate-1A abstract synchronized two-view fibre challenge
$N_\nu=X^{o(1)}$ together with unit-twist portability; the Gate-1B one-completion coefficient
$\beta_{D,P}=\mu_D\Lambda_P$, fixed-unit reciprocal-source invariance, the exact $1/104$ dyadic margin,
and the open source-to-convolution-BV interface; and the balanced-$R_9$ finite binomial certificate.
The physical Ford $G(d;n)$ specialization, the generated-(7.23) census, the replacement certificate
$C_{\rm FM}$, the Gate-1A necessity decision and all Gate-2 analytic inputs remain open.

\subsection*{Historical technical provenance}
The following remain useful diagnostics but are not part of the controlling public proof
architecture:
\begin{itemize}[leftmargin=1.6em,itemsep=0.35em]
\item the squarefree polarization plus the $X^{6/7+o(1)}$ repeated-prime router;
\item the V14 Vaughan/prime-modulus M\"obius two-outer route;
\item failed direct Gate-1A and native-QK dictionaries;
\item separate multiplicative-large-sieve and equal-seven three-block death tests.
\end{itemize}
They belong in the technical verification dossier rather than in the main narrative.

\begin{thebibliography}{9}
\bibitem{ErdosGraham1980}
P. Erd\H{o}s and R. L. Graham,
\emph{Old and New Problems and Results in Combinatorial Number Theory},
Monographies de L'Enseignement Math\'ematique, 1980, p.~33.

\bibitem{ErdosProblems287}
T. F. Bloom,
\emph{Erd\H{o}s Problem \#287}, Erd\H{o}s Problems archive,
problem provenance and status page, accessed August 2026.

\bibitem{FordMaynard2024}
K. Ford and J. Maynard,
\emph{On the theory of prime producing sieves},
arXiv:2407.14368, 2024.

\bibitem{DGS2017}
S. Drappeau, A. Granville and X. Shao,
\emph{Smooth-supported multiplicative functions in arithmetic progressions beyond the
$x^{1/2}$-barrier}, Mathematika 63 (2017), 895--918.

\bibitem{Pascadi2025}
A. Pascadi,
\emph{On the exponents of distribution of primes and smooth numbers},
arXiv:2505.00653v2, 29 June 2025.

\bibitem{LeanMathlib}
The Lean and Mathlib communities,
\emph{Lean 4 and Mathlib}, formal proof assistant and mathematical library.
\end{thebibliography}

\end{document}
